\chapter{Preliminaries: Set theory and categories}
\label{chap:prelims}

\emph{Set theory} is a mathematical field in itself, and its proper
treatment (say via the famous `\name{Zermelo}-\name{Fr\"ankel}'
axioms) goes well beyond the scope of this book and the competence of
this writer. We will only deal with so-called `naive' set theory,
which is little more than a system of notation and terminology
enabling us to express precisely mathematical definitions, statements,
and their proofs.

Familiarity with this language is essential in approaching a
subject such as algebra, and indeed the reader is assumed to have been
previously exposed to it.  In this chapter we first review some of the
language of naive set theory, mainly in order to establish the
notation we will use in the rest of the book. We will then get a small
taste of the language of categories, which plays a powerful unifying
role in algebra and many other fields. Our main objective is to convey
the notion of `universal property', which will be a constant refrain
throughout this book.

\section{Naive set theory}
\label{sec:prelims:naive-set-theory}

\subsection{Sets}
\label{subsec:prelims:sets}
The notion of set formalizes the intuitive idea of `collection of
objects'.  A set is determined by the \emph{elements} it contains: two
sets $A, B$ are equal (written $A = B$) if and only if they contain
precisely the same elements. `What is an \emph{element}?' is a
forbidden question in naive set theory: the buck must stop somewhere.
We can conveniently pretend that a `universe' of elements is available
to us, and we draw from this universe to construct the elements and
sets we need, implicitly assuming that all the operations we will
explore can be performed within this universe. (This is the tricky
point!) In any case, we specify a set by giving a precise recipe
determining which elements are in it. This definition is usually put
between braces and may consist of a simple, complete, list of
elements:
\[A \df {1, 2, 3}\] is\footnote{$\df$ is a notation often used to mean
  that the symbol on the left-hand side is defined by whatever is on
  the right-hand side. Logically, this is just expressing the equality
  of the two sides and could just as well be written `='; the extra :
  is a psychologically convenient decoration inherited from computer
  science.} the set consisting of the integers $1$, $2$, and $3$. By
convention, the order\footnote{Ordered lists are denoted with round
  parentheses: $(1, 2, 3)$ is not the same as $(1, 3, 2)$.} in which
the elements are listed, or repetitions in the list, are immaterial to
the definition.  Thus, the same set may be written out in many ways:
${1, 2, 3} = {1, 3, 2} = {1, 2, 1, 3, 3, 2, 3, 1, 1, 2, 1, 3}$.  This
way of denoting sets may be quite cumbersome and in any case will only
really work for \emph{finite sets}. For infinite sets, a popular way
around this problem is to write a list in which some of the elements
are understood as being part of a pattern---for example, the set of
even integers may be written $E = {\dots , -2, 0, 2, 4, 6, \dots }$,
but such a definition is inherently ambiguous, so this leaves room for
misinterpretation. Further, some sets are simply `too big' to be
listed, even in principle: for example (as one hopefully learns in
advanced calculus) there are simply too many real numbers to be able
to `list' them as one may `list' the integers.

It is often better to adopt definitions that express the elements of a
set as elements $s$ of some larger (and already known) set $S$
satisfying some property $P$.  One may then write
$A = \{s \in S \mid s\ \mathrm{satisfies}\ P\}$
($\in$ means element of\dots) and
this is in general precise and unambiguous\footnote{But note that
  there exist pathologies such as \name{Russell}'s paradox, showing that even
  this style of definitions can lead to nonsense. All is well so long
  as $S$ is indeed known to be a set to begin with.}.

We will
occasionally encounter a variation on the notion of set, called
`multiset'.  A multiset is a set in which the elements are allowed to
appear `with multiplicity': that is, a notion for which $\{2, 2\}$ would
be distinct from $\{2\}$. The correct way to define a multiset is by
means of functions, which we will encounter soon (see CHECKExample
2.2).  A few famous sets are
\begin{itemize}
\item $\emptyset$ : the empty set, containing no elements;
\item $\N$: the set of natural numbers (that is, nonnegative integers);
\item $\Z$: the set of integers;
\item $\Q$: the set of rational numbers;
\item $\R$: the set of real numbers;
\item $\C$: the set of complex numbers.
\end{itemize}
Also, the term \emph{singleton} is used to refer to any set consisting
of precisely one element. Thus ${1}$, ${2}$, ${3}$ are different sets,
but they are all singletons.

Here are a few useful symbols (called
quantifiers):
\begin{itemize}
\item $\exists$  means \emph{there exists}\ldots (the existential quantifier);
\item  $\forall$  means \emph{for all}\ldots (the universal quantifier).
\end{itemize}
Also, $\exists !$ is used to mean there exists a unique\ldots

For example, the set of even integers may be written as
\[E = \{a \in Z | (\exists n \in Z) a = 2n\}:\] in words, ``all
integers $a$ such that there exists an integer $n$ for which $a =
2n$''. In this case we could replace $\exists$ by $\exists !$ without
changing the set---but that has to do with properties of $Z$, not with
mathematical syntax. Also, it is common to adopt the shorthand
\[E = \{2n | n \in Z\], in which the existential quantifier is
understood.  Being able to parse such strings of symbols effortlessly,
and being able to write them out fluently, is extremely important. The
reader of this book is assumed to have already acquired this skill.

Note that the \emph{order} in which things are written may make a big
difference. For example, the statement
\[(\forall a \in  Z) (\exists b \in  Z)\, b = 2a\]
is true: it says that the result of doubling an arbitrary integer yields an
integer; but
\[(\exists b \in Z) (\forall a \in Z) b = 2a\] is false: it says that
there exists a fixed integer $b$ which is `simultaneously' twice as
much as every integer---there is no such thing.

Note also that writing simply
\[b = 2a\] by itself does not convey enough information, unless the
context makes it completely clear what quantifiers are attached to $a$
and $b$: indeed, as we have just seen, different quantifiers may make
this into a true or a false statement.

\subsection{Inclusion of sets}
\label{subsec:prelims:inclusion-of-sets}
As mentioned above, two sets are equal if and only if they contain the
same elements. We say that a set $S$ is a subset of a set $T$ if every
element of $S$ is an element of $T$ , in symbols, $S \subseteq T$.  By
convention, $S \subset T$ means the same thing: that is (unlike $<$
vs. $\leq$ ), it does not exclude the possibility that $S$ and $T$ may
be equal. To avoid any confusion, I will consistently use $\subseteq$ in
this book. One adopts $S \subsetneq T$ to mean that $S$ is `properly'
contained in $T$ : that is, $S \subseteq T$ and $S \ne T$.

We can think of `inclusion of sets' in terms of logic: $S \subseteq T$
means that $s \in S \Rightarrow s \in T$ (the quantifier $\forall s$ is
understood); that is, `if $s$ is an element of $S$, then $s$ is an
element of $T$'; that is, all elements of $S$ are elements of $T$;
that is, $S \subseteq T$ as promised.

Note that for all sets $S$, $\emptyset \subseteq S$ and
$S \subseteq S$.  If $S \subseteq T$ and $T \subseteq S$, then
$S = T$.  The symbol $|S|$ denotes the number of elements of $S$, if
this number is finite; otherwise, one writes $|S| = \infty$. If $S$
and $T$ are finite, then $S \subseteq T \Rightarrow |S| \leq |T |$.  The
subsets of a set $S$ form a set, called the power set, or the set of
parts of $S$. For example, the power set of the empty set $\emptyset$
consists of one element: $\emptyset$. The power set of $S$ is denoted
$\mathcal{P}(S)$; a popular alternative is $2^S$, and indeed
$|P(S)| = 2^{|S|}$ if $S$ is finite (cf. CHECKExercise 2.11).

\subsection{Operations between sets}
\label{subsec:prelims:operations-between-sets}
Once we have a few sets to play with, we can
obtain more by applying certain standard operations. Here are a few:
\begin{itemize}
\item $\cup$ : the union;
\item $\cap$ : the intersection;
\item $\setminus$ : the difference;
\item $\amalg$ : the disjoint union;
\item $\times$ : the (Cartesian) product; 
\item  and the important notion of `quotient by an equivalence relation'.
\end{itemize}

Most of these operations should be familiar to the reader: for example,
$\{1, 2, 4\} \cup  \{3, 4, 5\} = \{1, 2, 3, 4, 5\}$
while
$\{1, 2, 4\} \setminus  \{3, 4, 5\} = \{1, 2\}$.
TODO: Venn Diagrams
% In terms of \name{Venn} diagrams of infamous `new math' memory:
% S
%  T
% S \cup  T
%  S \cap  T
%  S \setminus  T
% (the solid black contour indicates the set included in the operation).
Several of these operations may be written out in a transparent way in terms
of logic: thus, for example,
\[
  s \in  S \cap  T \iff  (s \in  S\ \mathrm{and\ } s \in  T ).
\]



Two sets $S$ and $T$ are disjoint if $S \cap T = \emptyset$, that is, if
no element is `simultaneously' in both of them.

The \emph{complement} of a subset $T$ in a set $S$ is the difference
set $S \setminus T$ consisting of all elements of $S$ which are not
in$T$. Thus, for example, the complement of the set of even integers
in $\Z$ is the set of odd integers.

The operations $\amalg$ , $\times$ , and quotients by equivalence
relations are slightly more mysterious, and it is very instructive to
contemplate them carefully. We will look at them in a particularly
naive way first and come back to them in a short while when we have
acquired more language and can view them from a more sophisticated
viewpoint.

\subsection{Disjoint unions, products}
\label{subsec:prelims:disjoint-unions-products}
One problem with these operations is that their output may not be defined as a
set, but rather as a set up to isomorphisms of sets, that is, up to bijections.
To make sense out of this, we have to talk about functions, and we will do that
in a moment.

Roughly speaking, the disjoint union of two sets $S$ and $T$ is a set
$S \amalg T$ obtained by first producing `copies' $S'$ and $T'$ of the sets $S$
and $T$, with the property that $S' \cap T' = \emptyset$, and then taking the
(ordinary) union of $S'$ and $T'$. The careful reader will feel uneasy, since
this `recipe' does not define one set: whatever it means to produce a `copy' of
a set, surely there are many ways to do so. This ambiguity will be clarified
below.

Nevertheless, note that we can say something about $S \amalg T$ even on these
very shaky grounds: for example, if $S$ consists of 3 elements and $T$ consists
of 4 elements, the reader should expect (correctly) that $S \amalg T$ consists
of 7 elements.

Products are marred by the same kind of ambiguity, but fortunately there is
a convenient convention that allows us to write down `one' set representing
the product of two sets $S$ and $T$: given $S$ and $T$, we let $S \times T$
be the set whose elements are the ordered
pairs\footnote{One can define the ordered pair $(s, t)$ as a set by setting
$(s, t) = \{s, \{s, t\}\}$: this carries the information of the elements
$s$, $t$, as well as conveying the fact that $s$ is special (= the first
element of the pair).} $(s, t)$ of elements of $S$ and $T$:
\[
  S \times T := \{(s, t) \mid s \in S,\, t \in T\}.
\]
Thus, if $S = \{1, 2, 3\}$ and $T = \{3, 4\}$, then
$S \times T = \{(1, 3), (1, 4), (2, 3), (2, 4), (3, 3), (3, 4)\}$.
For a more sophisticated example, $\R \times \R$ is the set of pairs of real
numbers, which (as we learn in calculus) is a good way to represent a plane.
The set $\Z \times \Z$ could be represented by considering the points in this
plane that happen to have integer coordinates. Incidentally, it is common to
denote these sets $\R^2$, $\Z^2$; and similarly, the product $A \times A$ of
a set by itself is often denoted $A^2$.

If $S$ and $T$ are finite sets, clearly $|S \times T| = |S||T|$.

Also note that we can use products to obtain explicit `copies' of sets as
needed for the disjoint union: for example, we could let
$S' = \{0\} \times S$, $T' = \{1\} \times T$, guaranteeing that $S'$ and
$T'$ are disjoint (why?); and there is an evident way to `identify' $S$ and
$S'$, $T$ and $T'$. Again, making this precise requires a little more
vocabulary.
The operations $\cup$, $\cap$, $\amalg$, $\times$ extend to operations on
whole `families' of sets: for example, if $S_1, \dots, S_n$ are sets, we
write
\[
  \bigcap_{i=1}^{n} S_i = S_1 \cap S_2 \cap \dots \cap S_n
\]
for the set whose elements are those elements which are simultaneously
elements of all sets $S_1, \dots, S_n$; and similarly for the other
operations. But note that while it is clear from the definitions that, for
example,
\[
  S_1 \cup S_2 \cup S_3
  = (S_1 \cup S_2) \cup S_3
  = S_1 \cup (S_2 \cup S_3),
\]
it is not so clear in what sense the sets
\[
  S_1 \times S_2 \times S_3, \quad
  (S_1 \times S_2) \times S_3, \quad
  S_1 \times (S_2 \times S_3)
\]
should be `identified' (where we can define the leftmost set as the set of
`ordered triples' of elements of $S_1$, $S_2$, $S_3$, by analogy with the
definition for two sets). In fact, again, we can really make sense of such
statements only after we acquire the language of functions. However, all such
statements do turn out to be true, as the reader probably expects; by virtue
of this fortunate circumstance, we can be somewhat cavalier and gloss over
such subtleties.
More generally, if $\mathcal{S}$ is a set of sets, we may consider sets
\[
  \bigcup_{S \in \mathcal{S}} S, \quad
  \bigcap_{S \in \mathcal{S}} S, \quad
  \coprod_{S \in \mathcal{S}} S, \quad
  \prod_{S \in \mathcal{S}} S
\]
for the union, intersection, disjoint union, product of all sets in
$\mathcal{S}$. There are important subtleties concerning these definitions:
for example, if all $S \in \mathcal{S}$ are nonempty, does it follow that
$\prod_{S \in \mathcal{S}} S$ is nonempty? The reader probably thinks so,
but (if $\mathcal{S}$ is infinite) this is a rather thorny issue, amounting
to the axiom of choice.

By and large, such subtleties do not affect the material in this course; we
will partly come to terms with them in due
time\footnote{The reader will have to employ the axiom of choice in some
exercises every now and then, even before we come back to these issues, but
this will probably happen below the awareness level, and so it should.},
when they become more relevant to the issues at hand (cf.~\S{}V.3).

\subsection{Equivalence relations, partitions, quotients}
\label{subsec:prelims:equivalence-relations}

Intuitively, a relation on elements of a set $S$ is some special affinity
among selections of elements of $S$. For example, the relation $<$ on the
set $\Z$ is a way to compare the size of two integers: since $2 < 5$, $2$
 `is related to' $5$ in this sense, while $5$ is not related to $2$ in the
same sense. For all practical purposes, what a relation `means' is completely captured
by which elements are related to which elements in the set. We would really
know all there is to know about $<$ on $\Z$ if we had a complete list of
all pairs $(a, b)$ of integers such that $a < b$. For example, $(2, 5)$ is
such a pair, while $(5, 2)$ is not.
This leads to a completely straightforward definition of the notion of
relation: a relation on a set $S$ is simply a subset $R$ of the product
$S \times S$. If $(a, b) \in R$, we say that $a$ and $b$ are `related by
$R$' and write
\[
  a \mathrel{R} b.
\]
Often we use fancier symbols for relations, such as $<$, $\leq$, $=$,
$\sim$, \dots.
The prototype of a well-behaved relation is `=', which corresponds to the
`diagonal'
\[
  \{(a, b) \in S \times S \mid a = b\}
  = \{(a, a) \mid a \in S\}
  \subseteq S \times S.
\]
Three properties of this very special relation turn out to be particularly
important: if $\sim$ denotes for a moment the relation $=$ of equality,
then $\sim$ satisfies
\begin{itemize}
\item reflexivity: $(\forall a \in S)\; a \sim a$;
\item symmetry: $(\forall a \in S)\; (\forall b \in S)\; a \sim b
  \Rightarrow b \sim a$;
\item transitivity: $(\forall a \in S)\; (\forall b \in S)\;
  (\forall c \in S),\; (a \sim b \text{ and } b \sim c)
  \Rightarrow a \sim c$.
\end{itemize}
That is, every $a$ is equal to itself; if $a$ is equal to $b$, then $b$
is equal to $a$; etc.
\begin{defn}{Equivalence relation}{def:prelims:equivalence-relation}
An \emph{equivalence relation} on a set $S$ is any relation $\sim$
satisfying these three properties.
\end{defn}

In terms of the corresponding subset $R$ of $S \times S$, `reflexivity'
says that the diagonal is contained in $R$; `symmetry' says that $R$ is
unchanged if flipped about the diagonal (that is, if every $(a, b)$ is
interchanged with $(b, a)$); while unfortunately `transitivity' does not
have a similarly nice pictorial translation.

The datum of an equivalence relation on $S$ turns out to be equivalent to
a type of information which looks a little different at first, that is, a
partition of $S$. A partition of $S$ is a family of disjoint nonempty
subsets of $S$, whose union is $S$: for example,
\[\mathcal{P} = \{\{1, 4, 7\}, \{2, 5, 8\}, \{3, 6\}, \{9\}\}\]
is a partition of the set $\{1, 2, 3, 4, 5, 6, 7, 8, 9\}$.

Here is how to get a partition of $S$ from a relation $\sim$ on $S$: for
every element $a \in S$, the equivalence class of $a$ (w.r.t.\ $\sim$) is
the subset of $S$ defined by
\[ [a]_\sim \df \{b \in S \mid b \sim a\};\]
then the equivalence classes form a partition $P_\sim$ of $S$ (Exercise 1.2).

Conversely (Exercise 1.3) every partition $P$ is the partition corresponding in
this fashion to an equivalence relation. Therefore, the notions of `equivalence
relation on $S$' and `partition of $S$' are really equivalent.

Now we can view $P_\sim$ as a set (whose elements are the equivalence
classes with respect to $\sim$). This is the quotient operation mentioned
in \cref{subsec:prelims:operations-between-sets}.

\begin{defn}{Quotient of set}{def:prelims:quotient-set}
The \emph{quotient} of the set $S$ with respect to the equivalence
relation $\sim$ is the set
\[\quot{S}{\sim} \df P_\sim\]
of equivalence classes of elements of $S$ with respect to $\sim$.
\end{defn}

\begin{exam}{}{}
  Take $S = \Z$, and let $\sim$ be the relation defined by
  \[a \sim b \iff a - b \text{ is even.} \] Then $\quot{\Z}{\sim}$ consists of
  two equivalence classes: $\quot{\Z}{\sim} = \{[0]_\sim, [1]_\sim\}$.
  Indeed, every integer $b$ is either even (and hence $b-0$ is even,
  so $b \sim 0$, and $b \in [0]_\sim$) or odd (and hence $b-1$ is
  even, so $b \sim 1$, and $b \in [1]_\sim$). This is of course the
  starting point of modular arithmetic, which we will cover in due
  detail later on (\S{}II.2.3).
\end{exam}

One way to think about this operation is that the equivalence relation
`becomes equality in the quotient': that is, two elements of the
quotient $\quot{S}{\sim}$ are equal if and only if the corresponding elements
in $S$ are related by $\sim$. In other words, taking a quotient is a
way to turn any equivalence relation into an equality. This
observation will be further formalized in `categorical terms' in a
short while (\S{}5.3).

\subsection*{Exercises}

Exercises marked with a \exerused are referred to from the text; exercises marked with a
\lnot  are referred to from other exercises. These referring exercises
       and sections are
listed in brackets following the current exercise; see the introduction for
further
clarifications, if necessary.
\begin{exer}{}{}
Locate a discussion of Russell's paradox, and understand it.
\end{exer}{}{}
\begin{exer}{}{} Prove that if $\sim$ is an equivalence relation on a
  set $S$, then the corresponding family $P_\sim$ defined in
  \cref{subsec:prelims:equivalence-relations} is indeed a partition of
  $S$: that is, its elements are nonempty, disjoint, and their union
  is $S$. [\S{}1.5]
\end{exer}
\begin{exer}{}{}
  \exerused Given a partition $P$ on a set $S$, show how to define a
  relation $\sim$ on $S$ such that $P$ is the corresponding
  partition. [\S{}1.5]
\end{exer}
\begin{exer}{}{}
  How many different equivalence relations may be defined on the set $\{1, 2,
   3\}$?
\end{exer}
\begin{exer}{}{}
  Give an example of a relation that is reflexive and symmetric but not
  transitive.
  What happens if you attempt to use this relation to define a partition on the
   set?
  Hint: Thinking about the second question will help you answer the first one.
\end{exer}
\begin{exer}{}{}
  Define a relation $\sim$ on the set $\R$ of real numbers by setting
$a \sim b \iff b-a \in \Z$. Prove that this is an equivalence
relation, and find a `compelling' description for
$\quot{\R}{\sim}$. Do the same for the relation $\approx$ on the plane
$\R \times \R$ defined by declaring $(a_1, a_2) \approx (b_1 , b_2)
\iff b_1 - a_1 \in \Z \text{ and } b_2 - a_2 \in \Z$.  [\S{}II.8.1,
II.8.10]
\end{exer}

\section{Functions between sets}
\label{sec:prelims:functions-between-sets}

\subsection{Definition}
\label{subsec:prelims:functions:definition}
A common thread we will follow for just about every structure
introduced in this book will be to try to understand both the type of structures
and the ways in which different instances of a given structure may interact.

Sets interact with each other through functions. It is tempting to think
of a function $f$ from a set $A$ to a set $B$ in `dynamic' terms, as a
way to `go from $A$ to $B$'. Similarly to the business with relations,
it is straightforward to formalize this notion in ways that do not need
to invoke any deep `meaning' of any given $f$: everything that can be
known about a function $f$ is captured by the information of which
element $b$ of $B$ is the image of any given element $a$ of $A$. This
information is nothing but a subset of $A \times B$:
\[\Gamma_f := \{(a, b) \in A \times B \mid b = f(a)\}
  \subseteq A \times B.\]
This set $\Gamma_f$ is the graph of $f$; officially, a function really `is' its
graph.\footnote{To be precise, it is the graph $\Gamma_f$
 together with the information of the source $A$ and the
 target $B$ of $f$. These are part of the data of the function.}

Not all subsets $\Gamma \subseteq A \times B$ correspond to (`are')
functions: we need to put one requirement on the graphs of functions,
which can be expressed as follows:
\[(\forall a \in A) (\exists !b \in B)\quad (a, b) \in \Gamma_f,\]
or (`in functional notation')
\[(\forall a \in A) (\exists !b \in B)\quad f(a) = b.\]
That is, a function must send each element $a$ of $A$ to exactly one
element of $B$, depending on $a$. `Multivalued functions' such as
$\pm \sqrt{x}$ (which are very important in, e.g., the study of
\name{Riemann} surfaces) are not functions in this sense.

To announce that $f$ is a function from a set $A$ to a set $B$, one
writes $f : A \to B$ or draws the following picture (`diagram'):
\[A \overset{f}{\to} B.\]
The action of a function $f : A \to B$ on an element $a \in A$ is
sometimes indicated by a `decorated' arrow, as in
\[a \mapsto f(a).\]

The collection of all functions from a set $A$ to a set $B$ is itself
a set\footnote{This is another `operation among sets', not listed in
  \S{}1.3. Can you see why we use $B^A$ for this set?  (Cf. Exercise
  2.10.)}, denoted $B^A$. If we take seriously the notion that a
function is really the same thing as its graph, then we can view $B^A$
as a (special) subset of the power set of $A \times B$.

Every set $A$ comes equipped with a very special function, whose graph is the
diagonal in $A \times A$: the identity function on $A$
\[\id_A : A \to A\]
defined by $(\forall a \in A) \id_A(a) = a$. More generally, the
inclusion of any subset $S$ of a set $A$ determines a function
$S \to A$, simply sending every element $s$ of $S$ to `itself' in $A$.

If $S$ is a subset of $A$, we denote by $f (S)$ the subset of $B$ defined by
\[f(S) := \{b \in B \mid (\exists a \in S)\ b = f(a)\}.\] That is,
$f(S)$ is the subset of $B$ consisting of all elements that are images
of elements of $S$ by the function $f$. The largest such subset,
$f(A)$, is called the image of $f$, denoted $\image f$.

Also, $f\big|_S$ denotes the `restriction' of $f$ to the subset $S$:
this is the function $S \to B$ defined by
\[(\forall s \in S) : f\big|_S (s) = f (s).\] That is, $f\big|_S$ is
the composition (in the sense explained in \Cref{XX}) $f \circ i$, where
$i : S \to A$ is the inclusion. Note that $f (S) = \image(f\big|_S)$.

\subsection{Examples: Multisets, indexed sets}
\label{subsec:prelims:examples-multisets-indexed-sets}
The `multisets' mentioned briefly in \S{}1.1 are a simple example of a
notion easily formalized by means of functions.  A multiset may be
defined by giving a function from a (regular) set $A$ to the set $\N^*$
of positive\footnote{Some references allow 0 as a possible
  multiplicity.} integers; if $m : A \to \N^*$ is such a function, the
corresponding multiset consists of the elements $a \in A$, each taken
$m(a)$ times. Thus, the multiset $\{a, a, a, b, b, b, b, b, c\}$ is
really the function $m : \{a, b, c\} \to \N^*$ for which $m(a) = 3$,
$m(b) = 5$, $m(c) = 1$. As with ordinary sets, the order in which the
elements are listed is not part of the information carried by a
multiset. Simple set-theoretic notions such as inclusion, union, etc.,
extend in a straightforward way to multisets.  For another viewpoint
on multisets, see \cref{XX}.

Another example is given by the use of `indices'. If we write let
$a_1,\dots ,a_n$ be integers \dots, we really mean consider a function
$a : \{1,\dots , n\} \to \Z\dots$, with the understanding that $a_i$
is shorthand for the value $a(i)$ (for $i = 1, \dots,n)$. It is
tempting to think of an indexed set $\{a_i\}_{i\in I}$ simply as a set
whose elements happen to be denoted $a_i$, for $i$ ranging over some
`set of indices' $I$; but such an indexed set is more properly a
function $I \to A$, where $A$ is some set from which we draw the
elements $a_i$. For example, this allows us to consider $a_0$ and
$a_1$ as distinct elements of $\{a_i\}_{i\in \N}$, even if by
coincidence $a_0 = a_1$ as elements of the target set $A$.

It is easy to miss such subtleties, and some abuse of notation is common and
usually harmless. These distinctions play a role in (for example) discussions of
linear independence of sets of vectors; cf. \S{}VI.1.2.

\subsection{Composition of functions}
\label{subsec:prelims:composition-of-functions}
Functions may be composed: if $f : A \to B$ and $g : B \to C$ are
functions, then so is the operation $g \circ f$ defined by
\begin{equation}
  \label{eq:comp-def}
  (\forall a \in A)\quad (g \circ f)(a) := g(f(a)) :
\end{equation}
that is, we use $f$ to go from $A$ to $B$, then apply $g$ to reach $C$.
Graphically we may draw pictures such as
\[
\begin{tikzcd}
  A \arrow[r, "f"] \arrow[rr, bend right, "g \circ f"'] & B \arrow[r, "g"] & C
\end{tikzcd}
\qquad
\begin{tikzcd}
  A \arrow[r,"f"] \arrow[rd,"g \circ f"'] & B \arrow[d,"g"]\\
  & C
\end{tikzcd}
\]
Such graphical representations of collections of (for example) sets
connected by functions are called \emph{diagrams}. We will draw many
diagrams, and in contexts substantially more general than the one at
hand right now.

We say that the diagrams drawn above `commute', or `are commutative',
meaning that if we start from $A$ and travel to $C$ in either of the
two possible ways prescribed by the diagram, the result of applying
the functions one encounters is the same. This is precisely the
content of the statement \eqref{eq:comp-def}.

Composition is associative: that is, if $f : A \to B$, $g : B \to C$,
and $h : C \to D$ are functions, then
$h \circ (g \circ f) = (h \circ g) \circ f$. Graphically, the diagram
\[
  \begin{tikzcd}
    A \arrow[r,"f"] \arrow[rr,bend right, "g\circ f" description]
    & B \arrow[r,"g"] \arrow[rr,bend left, "h\circ g" description]
    & C \arrow[r,"h"]
    & D
  \end{tikzcd}
\]
commutes. This important observation should be completely evident from the
definition of composition.

The identity function is very special with respect to compositions: if
$f : A \to B$ is any function, then $\id_B \circ f = f$ and
$f \circ \id_A = f$. Graphically, the diagrams
\[
  \begin{tikzcd}
    A \arrow[r,"f"] \arrow[rr,bend right,"f"']& B \arrow[r,"\id_B"] & B
  \end{tikzcd}
  \qquad
  \begin{tikzcd}
    A \arrow[r,"\id_A"] \arrow[rr,bend right,"f"']& A \arrow[r,"f"]& B
  \end{tikzcd}
\]
commute.

\subsection{Injections, surjections, bijections}
\label{subsec:prelims:injections-surjections-bijections}
Special kinds of functions deserve highlighting:
\begin{enumerate}
\item A function $f : A \to B$ is injective (or an injection or
  one-to-one) if
  \[\kforall_{a' \in A}\ \kforall_{a'' \in A}\
    a' = a'' \Rightarrow f(a') = f(a'')\]
  that is, if $f$ sends different elements to different elements
  \footnote{Often one checks this definition in the contrapositive
    (hence equivalent) formulation, that is,
    $(\forall a' \in A)(\forall a'' \in A)
    \quad f(a') = f(a'') \Rightarrow a' = a''$.}.
\item A function $f : A \to B$ is surjective (or a surjection or onto)
  if
  \[\kforall_{b \in B}\ \kexists_{a \in A}\ b = f(a):\]
  that is, if $f$ `covers the whole of $B$'; more precisely, if
  $\image f = B$.
\end{enumerate}

Injections are often drawn $\hookrightarrow$; surjections are often drawn
$\twoheadrightarrow$.

If $f$ is both injective and surjective, we say it is bijective (or a
bijection or a one-to-one correspondence or an isomorphism of sets.)
In this case we often write $f : A \xrightarrow{\sim} B$, or
\[A \overset{\sim}{=} B,\]
and we say that $A$ and $B$ are `isomorphic' sets.

Of course the identity function $\id_A : A \to A$ is a bijection.

If $A \overset{\sim}{=} B$, that is, if there is a bijection
$f : A \to B$, then the sets $A$ and $B$ may be `identified' through
$f$, in the sense that we can match precisely the elements $a$ of $A$
with the corresponding elements $f(a)$ of $B$. For example, if $A$ is
a finite set and $A \overset{\sim}{=} B$, then $B$ is necessarily
also a finite set and $|A| = |B|$.

This terminology allows us to make better sense of the considerations
on `disjoint union' given in \S{}1.3: the `copies' $A'$, $B'$ of the
given sets $A$, $B$ should simply be isomorphic sets to $A$, $B$,
respectively. The proposal given at
the end of \S{}1.4 to produce such disjoint `copies' works, because
(for example) the function $f : A \to \{0\} \times A$ defined by
\[\kforall_{a \in A} f (a) = (0, a)\] is manifestly a bijection.

\subsection{Injections, surjections, bijections: Second viewpoint}
\label{subsec:prelims:injections-surjections-bijections-second-viewpoint}
There is an alternative and instructive way to think about these notions.
If $f : A \to  B$ is a bijection, then we can `flip its graph' and define a
function
$g : B \to  A$ :
that is, we can let $a = g(b)$ precisely when $b = f (a)$. (The fact that $f$ is both
injective and surjective guarantees that the flip of $\Gamma_f$ is the graph of a
function
according to the definition given in \S{}2.1. Check this!)
This function $g$ has a very interesting property: graphically,
\[
  \begin{tikzcd}
    A \arrow[r,"f"] \arrow[rr,bend right,"\id_A" description]& B\arrow[r,"g"] &A
  \end{tikzcd}
  \quad
    \begin{tikzcd}
     B\arrow[r,"g"]\arrow[rr,bend right,"\id_B" description] &A\arrow[r,"f"]&B
  \end{tikzcd}
\] commute; that is, $g \circ f = \id_A$ and $f \circ g = \id_B$. The
first identity tells us that $g$ is a `left-inverse'\footnote{Never
  mind that $g$ is drawn to the right of $f$ in the diagram---we say
  that $g$ is a left-inverse of $f$ because it is written to the left
  of $f : g \circ f = \id_A$} of $f$; the second tells us that $g$ is
a `right-inverse' of $f$. We simply say that it is the inverse of $f$, denoted $f^{-1}$. Thus, `bijections have inverses'.

What about the converse? If a function has an inverse, is it a bijection? This
is true, but in fact we can be much more precise.

\begin{prop}{}{prop:inv-inj-sur}
  Assume $A \ne\emptyset$, and let $f : A \to  B$ be a function. Then
  \begin{itemize}
  \item $f$ has a left-inverse if and only if it is injective.
  \item $f$ has a right-inverse if and only if it is surjective.
  \end{itemize}
\end{prop}
\begin{proof}
  ( $\Rightarrow$  ) If $f : A \to  B$ has a left-inverse, then
  there exists a $g : B \to
A$ such that
$g \circ  f = id_A$ . Now assume that $a' \ne a''$  are
arbitrary different elements in $A$ then
\[  
  g(f(a') = \id_A(a') = a' \ne a'' = \id_A(a'') = g(f(a''))
\]           
that is, $g$ sends $f (a\setminus)$ and $f (a'' )$ to different
elements. This forces $f (a')$ and $f (a'')$
to be different, showing that $f$ is injective.

( $\Leftarrow$ ) Now assume $f : A \to  B$ is injective. In order to construct a
function
$g : B \to  A$, we have to assign a unique value $g(b) \in  A$ for each element  $b
\in  B$. For
this, choose any fixed element $s \in  A$ (which we can do because $A \ne \emptyset$
); then set
\[
  g(b) :=
  \begin{cases}
    a & \text{if } b = f(a) \text{ for some } a \in A,\\
    s & \text{if } b \not\in \image f.
  \end{cases}
\]
In words, if $b$ is the image of an element $a$ of $A$, send it back
to $a$; otherwise, send it to your fixed element $s$.

The given assignment defines a function, precisely because $f$ is
injective: indeed, this guarantees that every $b$ that is the image of
some $a \in A$ by $f$ is the image of a unique a (two distinct
elements of A cannot be simultaneously sent to $b$ by $f$, since $f$
is injective). Thus every $b \in B$ is sent to a unique well-defined
element of $A$, as is required of functions.  Finally, the function
$g : B \to A$ is a left-inverse of $f$ . Indeed, if $a \in A$, then
$b = f (a)$ is of the first type, so it is sent back to a by $g$; that
is, $g \circ f(a) = a = \id_A(a)$ for all $a \in A$, as needed.  The
proof of (2) is left as an exercise (Exercise 2.2).
\end{proof}

\begin{coro}{}{}
  A function $f : A \to  B$ is a bijection if and only if it has a
  (two-sided) inverse.
\end{coro}

This is not completely innocent: if $f$ has both a left-inverse and a
right-inverse, why should it have one inverse that works as both on
the left and on the right?  Try to prove this by yourself now. We will
come back to this issue soon (in \S{}4).

If a function is injective but not surjective, then it will not have a
right-inverse, and if the source has at least two elements, it will
necessarily have more than one left-inverse (this should be clear from
the argument given in the proof of Proposition 2.1). Similarly, a
surjective function will in general have many right-inverses; they are
often called sections.

\Cref{prop:inv-inj-sur} hints that something deep is going on
here. The definition of injective and surjective maps given in \S{}2.4
relied crucially on working directly with the elements of our sets;
\Cref{prop:inv-inj-sur} shows that in fact these properties are
detected by the way functions are `organized' among sets. Even if we
did not know what `elements' means, still we could make sense of the
notions of injectivity and surjectivity (and hence of isomorphisms of
sets) by exclusively referring to properties of functions.

This is a more `mature' point of view and one that will be championed
when we talk about categories. To some extent, it should cure the
reader from the discomfort of talking about `elements', as we did in
our informal introduction to sets, without defining what these
mysterious entities are supposed to be.

The standard notation for \emph{the} inverse of a bijection $f$ is
$f^{-1}$. This symbol is also used for functions that are not
bijections, but in a slightly different context: if $f : A \to B$ is
any function and $T \subseteq B$ is a subset of $B$, then $f^{-1} (T)$
denotes the subset of $A$ of `all elements that map to $T$'; that is,
$f^{-1}(T) = \{a \in A \mid f (a) \in T \}$.  If $T = {q}$ consists of
a single element of $B$, $f^{-1}(T)$ (abbreviated $f^{-1}(q)$) is
called the \emph{fiber} of $f$ over $q$. Thus a function $f : A \to B$
is a bijection if it has nonempty fibers over all elements of $B$
(that is, $f$ is surjective), and these fibers are in fact singletons
(that is, $f$ is injective). In this case, this notation $f$ matches
nicely with the notation of `inverse' mentioned above.

\subsection{Monomorphisms and epimorphisms}
\label{subsec:prelims:functions:monomorphisms-and-epimorphisms}
There is yet another way to express injectivity and surjectivity, which
appears at first more complicated than what
we have seen so far but which is in fact even more basic.

A function $f : A \to  B$ is a monomorphism (or monic) if the following holds:
for all sets $Z$ and all functions $\alpha', \alpha'' Z \to  A$
\[f \circ  \alpha' = f \circ  \alpha''  \Rightarrow
  \alpha'  = \alpha''.\]
\begin{prop}{}{prop:old-2.3}
  A function is injective if and only if it is a monomorphism.
\end{prop}
\begin{proof}
 ( \Rightarrow  ) By Proposition 2.1 \Cref{XX}, if a function $f : A \to  B$ is
injective, then it
has a left-inverse $g : B \to  A$. Now assume that $\alpha'  ,\alpha''$  are arbitrary functions from
another set $Z$ to $A$ and that
\[f \circ  \alpha'  = f \circ  \alpha'' ;\]
compose on the left by $g$, and use associativity of composition:
\[(g \circ  f ) \circ  \alpha'  = g \circ  (f \circ  \alpha')
= g \circ  (f \circ  \alpha''  ) = (g \circ  f ) \circ  \alpha'' ;\] since $g$ is a left-inverse of $f$ , this says
\[\id_A \circ  \alpha'  = \id_A \circ  \alpha'',\]
and therefore
\[\alpha'  = \alpha'',\]
as needed to conclude that $f$ is a monomorphism.

( \Leftarrow = ) Now assume that $f$ is a monomorphism. This says something about
arbitrary sets $Z$ and arbitrary functions $Z \to  A$; we are going to use a
microscopic
portion of this information, choosing $Z$ to be any singleton ${p}$. Then assigning
functions $\alpha' , \alpha'' : Z \to  A$ amounts to
choosing to which elements $a' = \alpha ' (p)$,
$a''  = \alpha''(p)$ we should send the
single element $p$ of $Z$. For this particular choice
of $Z$, the property defining monomorphisms, $f \circ  \alpha'  = f \circ
\alpha''  \Rightarrow  \alpha '  = \alpha'$
  , becomes
\[f \circ  \alpha'(p) = f \circ  \alpha'' (p)
\Rightarrow  \alpha'  = \alpha''\],
that is,
$f (a') = f (a'') \Rightarrow  \alpha'  =
\alpha''$.

Now two functions from $Z = {p}$ to $A$ are equal if and only if they send $p$ to the
same element, so this says
\[f (a') = f (a'') \Rightarrow  a'  = a''  .\]
This has to be true for all $\alpha', \alpha''$,
that is, for all choices of distinct $a', a''$  in $A$.
In
other words, $f$ has to be injective, as was to be shown.
\end{proof}
The reader should now expect that there be a definition in the style
of the one given for monomorphisms and which will turn out to be
equivalent to `surjective'.  This is the case: such a notion is called
epimorphism. Finding it, and proving the equivalence with the ordinary
definition of `surjective', is left to the reader\footnote{This is a particularly important exercise, and I recommend that the reader
write out all
the gory details carefully.} (Exercise 2.5).


\subsection{Basic examples}
\label{subsec:prelims:basic-examples}
The basic operations on sets provided us with several important
examples of injective and surjective functions.

\begin{exam}{}{}
  Let $A$, $B$ be sets. Then there are natural \emph{projections} $\pi_A$, $\pi_B$:
  \[
    \begin{tikzcd}
      & A \times B \arrow[dl, "\pi_A"'] \arrow[dr, "\pi_B"] & \\
      A & & B
    \end{tikzcd}
  \]
  defined by
  \[ \pi_A((a,b)) := a, \quad \pi_B((a,b)) := b \]
  for all $(a,b) \in A \times B$. Both of these maps are (clearly) surjective.
\end{exam}

\begin{exam}{}{}
  Similarly, there are natural injections from $A$ and $B$ to the disjoint union:
  \[
    \begin{tikzcd}
      A \arrow[dr, hook] & & B \arrow[dl, hook'] \\
      & A \sqcup B &
    \end{tikzcd}
  \]
  obtained by sending $a \in A$ (resp., $b \in B$) to the corresponding element in
  the isomorphic copy $A'$ of $A$ (resp., $B'$ of $B$) in $A \sqcup B$.
\end{exam}

\begin{exam}{}{}
  If $\sim$ is an equivalence relation on a set $A$, there is a (clearly
  surjective) canonical projection
  \[ A \twoheadrightarrow A/{\sim} \]
  obtained by sending every $a \in A$ to its equivalence class $[a]_\sim$.
\end{exam}

\subsection{Canonical decomposition}
\label{subsec:prelims:canonical-decomposition}
The reason why we focus our attention on injective and surjective maps is that
they provide the basic `bricks' out of which any function may be constructed.

To see this, we observe that every function $f : A \to B$ determines an
equivalence relation $\sim$ on $A$ as follows: for all $a', a'' \in A$,
\[ a' \sim a'' \iff f(a') = f(a''). \]
(The reader should check that this is indeed an equivalence relation.)

\begin{thm}{}{thm:canonical-decomposition}
  Let $f : A \to B$ be any function, and define $\sim$ as above. Then $f$
  decomposes as follows:
  \[
    \begin{tikzcd}
      A \arrow[r, two heads] \arrow[rrr, bend left=20, "f"] &
      (A/{\sim}) \arrow[r, "\tilde{f}", "\sim"'] &
      \image f \arrow[r, hook] &
      B
    \end{tikzcd}
  \]
  where the first function is the canonical projection $A \to A/{\sim}$ (as in
  Example~2.6), the third function is the inclusion $\image f \subseteq B$, and
  the bijection $\tilde{f}$ in the middle is defined by
  \[ \tilde{f}([a]_\sim) := f(a) \]
  for all $a \in A$.
\end{thm}

The formula defining $\tilde{f}$ shows immediately that the diagram commutes; so
all we have to verify in order to prove this theorem is that
\begin{itemize}
\item that formula does define a function;
\item that function is in fact a bijection.
\end{itemize}
The first item is an instance of a class of verifications of the utmost
importance. The formula given for $\tilde{f}$ has a colossal built-in ambiguity:
the same element in $A/{\sim}$ may be the equivalence class of many elements of
$A$; applying the formula for $\tilde{f}$ requires choosing one of these elements
and applying $f$ to it. We have to prove that the result of this operation is
independent from this choice: that is, that all possible choices of
representatives for that equivalence class lead to the same result.

We encode this type of situation by saying that we have to verify that $\tilde{f}$
is \emph{well-defined}. We will often have to check that the operations we
consider are well-defined, in contexts very similar to the one epitomized here.

\begin{proof}
  Spelling out the first item discussed above, we have to verify that, for all
  $a', a''$ in $A$,
  \[ [a']_\sim = [a'']_\sim \implies f(a') = f(a''). \]
  Now $[a']_\sim = [a'']_\sim$ means that $a' \sim a''$, and the definition of
  $\sim$ has been engineered precisely so that this would mean $f(a') = f(a'')$
  as required here. So $\tilde{f}$ is indeed well-defined.

  To verify the second item, that $\tilde{f} : A/{\sim} \to \image f$ is a
  bijection, we check explicitly that $\tilde{f}$ is injective and surjective.

  \emph{Injective:} If $\tilde{f}([a']_\sim) = \tilde{f}([a'']_\sim)$, then
  $f(a') = f(a'')$ by definition of $\tilde{f}$; hence $a' \sim a''$ by
  definition of $\sim$, and then $[a']_\sim = [a'']_\sim$. Therefore
  \[ \tilde{f}([a']_\sim) = \tilde{f}([a'']_\sim) \implies [a']_\sim = [a'']_\sim, \]
  proving injectivity.

  \emph{Surjective:} Given any $b \in \image f$, there is an element $a \in A$
  such that $f(a) = b$. Then
  \[ \tilde{f}([a]_\sim) = f(a) = b \]
  by definition of $\tilde{f}$. Since $b$ was arbitrary in $\image f$, this shows
  that $\tilde{f}$ is surjective, as needed.
\end{proof}

\Cref{thm:canonical-decomposition} shows that every function is the composition
of a surjection, followed by an isomorphism, followed by an injection. While its
proof is trivial, this is a result of some importance, since it is the prototype
of a situation that will occur several times in this book. It will resurface every
now and then, with names such as `the first isomorphism theorem'.

\subsection{Clarification}
\label{subsec:prelims:clarification}

Finally, we can begin to clarify one comment about disjoint unions, products,
and quotients, made in \S{}1.4. Our definition of $A \sqcup B$ was the
(conventional) union of two disjoint sets $A'$, $B'$ isomorphic to $A$, $B$,
respectively. It is easy to provide a way to effectively produce such isomorphic
copies (as we did in \S{}1.4); but it is in fact a little too easy---many other
choices are possible, and one does not look any better than any other. It is in
fact more sensible not to make a fixed choice once and for all and simply accept
the fact that all of them produce acceptable candidates for $A \sqcup B$. From
this egalitarian standpoint, the result of the operation $A \sqcup B$ is not
`well-defined' as a set in the sense specified above. However, it is easy to see
(Exercise~2.9) that $A \sqcup B$ is well-defined up to isomorphism: that is,
that any two choices for the copies $A'$, $B'$ lead to isomorphic candidates for
$A \sqcup B$. The same considerations apply to products and quotients.

The main feature of sets obtained by taking disjoint unions, products, or
quotients is not really `what elements they contain' but rather `their
relationship with all other sets'. This will be (even) clearer when we revisit
these operations and others\footnote{The reader should also be aware that there
  are important variations on the operations we have seen so far---particularly
  important are the \emph{fibered} flavors of products and disjoint unions.} in
the context of categories.

\subsection*{Exercises}

\begin{exer}{}{}
  \exerused How many different bijections are there between a set $S$ with $n$
  elements and itself? [\S{}II.2.1]
\end{exer}

\begin{exer}{}{}
  \exerused Prove statement~(2) in \Cref{prop:inv-inj-sur}. You may assume that
  given a family of disjoint nonempty subsets of a set, there is a way to choose
  one element in each member of the family.\footnote{This (reasonable) statement
    is the \emph{axiom of choice}; cf.\ \S{}V.3.} [\S{}2.5, V.3.3]
\end{exer}

\begin{exer}{}{}
  Prove that the inverse of a bijection is a bijection and that the composition of
  two bijections is a bijection.
\end{exer}

\begin{exer}{}{}
  \exerused Prove that `isomorphism' is an equivalence relation (on any set of
  sets). [\S{}4.1]
\end{exer}

\begin{exer}{}{}
  \exerused Formulate a notion of \emph{epimorphism}, in the style of the notion
  of monomorphism seen in \S{}2.6, and prove a result analogous to
  \Cref{prop:old-2.3}, for epimorphisms and surjections. [\S{}2.6, \S{}4.2]
\end{exer}

\begin{exer}{}{}
  With notation as in Example~2.4, explain how any function $f : A \to B$
  determines a section of $\pi_A$.
\end{exer}

\begin{exer}{}{}
  Let $f : A \to B$ be any function. Prove that the graph $\Gamma_f$ of $f$ is
  isomorphic to $A$.
\end{exer}

\begin{exer}{}{}
  Describe as explicitly as you can all terms in the canonical decomposition
  (cf.\ \S{}2.8) of the function $\R \to \C$ defined by $r \mapsto e^{2\pi i r}$.
  (This exercise matches one assigned previously. Which one?)
\end{exer}

\begin{exer}{}{}
  \exerused Show that if $A' \cong A''$ and $B' \cong B''$, and further
  $A' \cap B' = \emptyset$ and $A'' \cap B'' = \emptyset$, then
  $A' \cup B' \cong A'' \cup B''$. Conclude that the operation $A \sqcup B$ (as
  described in \S{}1.4) is well-defined up to isomorphism (cf.\ \S{}2.9).
  [\S{}2.9, 5.7]
\end{exer}

\begin{exer}{}{}
  \exerused Show that if $A$ and $B$ are finite sets, then $|B^A| = |B|^{|A|}$.
  [\S{}2.1, 2.11, \S{}II.4.1]
\end{exer}

\begin{exer}{}{}
  \exerused In view of Exercise~2.10, it is not unreasonable to use $2^A$ to
  denote the set of functions from an arbitrary set $A$ to a set with 2 elements
  (say $\{0,1\}$). Prove that there is a bijection between $2^A$ and the power
  set of $A$ (cf.\ \S{}1.2). [\S{}1.2, III.2.3]
\end{exer}

\section{Categories}
\label{sec:prelims:categories}

The language of categories is affectionately known as \emph{abstract nonsense},
so named by \name{Norman Steenrod}. This term is essentially accurate and not
necessarily derogatory: categories refer to nonsense in the sense that they are
all about the `structure', and not about the `meaning', of what they represent.
The emphasis is less on how you run into a specific set you are looking at and
more on how that set may sit in relationship with all other sets. Worse (or
better) still, the emphasis is less on studying sets, and functions between sets,
than on studying `things, and things that go from things to things' without
necessarily being explicit about what these things are: they may be sets, or
groups, or rings, or vector spaces, or modules, or other objects that are so
exotic that the reader has no right whatsoever to know about them (yet).

`Categories' will intuitively look like sets at first, and in multiple ways.
Categories may make you think of sets, in that they are `collections of objects',
and further there will be notions of `functions from categories to categories'
(called \emph{functors}\footnote{However, we will not consider functors until
  later chapters: our first formal encounter with functors will be in Chapter
  VIII.}). At the same time, every category may make you think of the collection
of all sets, since there will be analogs of `functions' among the things it
contains.

\subsection{Definition}
\label{subsec:prelims:categories:definition}

The definition of a category looks complicated at first, but the gist of it may
be summarized quickly: a category consists of a collection of `objects', and of
`morphisms' between these objects, satisfying a list of natural conditions.

The reader will note that I refrained from writing a \emph{set} of objects,
opting for the more generic `collection'. This is an annoying, but unavoidable,
difficulty: for example, we want to have a `category of sets', in which the
`objects' are sets and the `morphisms' are functions between sets, and the
problem is that there simply is not a set of all
sets\footnote{That is one thing we learn from Russell's paradox.}. In a sense,
the collection of all sets is `too big' to be a set. There are however ways to
deal with such `collections', and the technical name for them is \emph{class}.
There is a `class' of all sets (and there will be classes taking care of groups,
rings, etc.).

An alternative would be to define a large enough set (called a \emph{universe})
and then agree that all objects of all categories will be chosen from this
gigantic entity.

In any case, all the reader needs to know about this is that there is a way to
make it work. We will use the term `class' in the definition, but this will not
affect any proof or any other definition in this book. Further, in some of the
examples considered below the class in question is a set (we say that the
category is \emph{small} in this case), so the reader will feel perfectly at
home when contemplating these examples.

\begin{defn}{}{defn:category}
  A \emph{category} $\mathsf{C}$ consists of
  \begin{itemize}
  \item a class $\operatorname{Obj}(\mathsf{C})$ of \emph{objects} of the
    category; and
  \item for every two objects $A$, $B$ of $\mathsf{C}$, a set
    $\operatorname{Hom}_{\mathsf{C}}(A, B)$ of \emph{morphisms},
  \end{itemize}
  with the properties listed below.
\end{defn}

As a prototype to keep in mind, think of the objects as `sets' and of morphisms
as `functions'. This one example should make the defining properties of morphisms
look natural and easy to remember:
\begin{itemize}
\item For every object $A$ of $\mathsf{C}$, there exists (at least) one morphism
  $1_A \in \operatorname{Hom}_{\mathsf{C}}(A, A)$, the `identity' on $A$.
\item One can compose morphisms: two morphisms
  $f \in \operatorname{Hom}_{\mathsf{C}}(A, B)$ and
  $g \in \operatorname{Hom}_{\mathsf{C}}(B, C)$ determine a morphism
  $gf \in \operatorname{Hom}_{\mathsf{C}}(A, C)$. That is, for every triple of
  objects $A$, $B$, $C$ of $\mathsf{C}$ there is a function (of sets)
  \[ \operatorname{Hom}_{\mathsf{C}}(A, B) \times \operatorname{Hom}_{\mathsf{C}}(B, C)
    \to \operatorname{Hom}_{\mathsf{C}}(A, C), \]
  and the image of the pair $(f, g)$ is denoted $gf$.
\item This `composition law' is associative: if $f \in \operatorname{Hom}_{\mathsf{C}}(A, B)$,
  $g \in \operatorname{Hom}_{\mathsf{C}}(B, C)$, and
  $h \in \operatorname{Hom}_{\mathsf{C}}(C, D)$, then $(hg)f = h(gf)$.
\item The identity morphisms are identities with respect to composition: that
  is, for all $f \in \operatorname{Hom}_{\mathsf{C}}(A, B)$ we have
  \[ f 1_A = f, \qquad 1_B f = f. \]
\end{itemize}

This is really a mouthful, but again, to remember all this, just think of
functions of sets. One further requirement is that the sets
$\operatorname{Hom}_{\mathsf{C}}(A, B)$ and $\operatorname{Hom}_{\mathsf{C}}(C, D)$
be disjoint unless $A = C$, $B = D$; this is something you do not usually think
about, but again it holds for ordinary set-functions.\footnote{I will often use
  the term `set-function' to emphasize that we are dealing with a function in
  the context of sets.} That is, if two functions are one and the same, then
necessarily they have the same source and the same target: source and target are
part of the datum of a set-function.

A morphism of an object $A$ of a category $\mathsf{C}$ to itself is called an
\emph{endomorphism}; $\operatorname{Hom}_{\mathsf{C}}(A, A)$ is denoted
$\operatorname{End}_{\mathsf{C}}(A)$. One of the axioms of a category tells us
that this is a `pointed' set, as $1_A \in \operatorname{End}_{\mathsf{C}}(A)$.
The reader should note that composition defines an `operation' on
$\operatorname{End}_{\mathsf{C}}(A)$: if $f$, $g$ are elements of
$\operatorname{End}_{\mathsf{C}}(A)$, so is their composition $gf$.

Writing `$f \in \operatorname{Hom}_{\mathsf{C}}(A, B)$' gets tiresome in the
long run. If the category is understood, one may safely drop the index
$\mathsf{C}$, or even use arrows as we do with set-functions: $f : A \to B$.
This also allows us to draw \emph{diagrams} of morphisms in any category; a
diagram is said to `commute' (or to be a `commutative' diagram) if all ways to
traverse it lead to the same results of composing morphisms along the way, just
as explained for diagrams of functions of sets in \S{}2.3.

In fact, I will now feel free to use diagrams as possible objects of categories.
The official definition of a diagram in this context would be a set of objects
of a category $\mathsf{C}$, along with prescribed morphisms between these
objects; the diagram commutes if it does in the sense specified above. The
specifics of the visual representation of a diagram are of course irrelevant.

\subsection{Examples}
\label{subsec:prelims:examples}
The reader should note that 90\% of the definition of the notion of category
goes into explaining the properties of its morphisms; it is fair to say that the
morphisms are the important constituents of a category. Nevertheless, it is
psychologically irresistible to think of a category in terms of its objects: for
example, one talks about the `category of sets'. The point is that usually the
kind of `morphisms' one may consider are (psychologically at least) determined
by the objects: if one is talking about sets, what can one possibly mean for
`morphism' other than a function of sets? In other situations (cf.\ Example~3.5
below or Exercise~3.9) it is a little less clear what the morphisms should be,
and looking for the `right' notion may be an interesting project.
\begin{exam}{}{}
  It is hopefully crystal clear by now that sets (as objects), together with
  set-functions (as morphisms), form a category; if not, the reader must stop
  here and go no further until this assertion sheds any residual
  mystery.\footnote{I will give the reader such prompts every now and then: at
    key times, it is more useful to take stock of what one knows than blindly
    march forward hoping for the best. A difficulty at this time signals the need
    to reread the previous material carefully. If the mystery persists, that's
    what office hours are there for. But typically you should be able to find
    your way out on your own, based on the information I have given you, and you
    will most likely learn more this way. You should give it your best try before
    seeking professional help.}

  There is no universally accepted, official notation for this important
  category. It is customary to write the word `Set' or `Sets', with some fancy
  decoration for emphasis. For example, in the literature one may encounter
  $\mathbf{Set}$, $\underline{\mathbf{Sets}}$, $\mathfrak{Set}$, $(Sets)$, and
  many amusing variations on these themes. We will use `sans-serif' fonts to
  denote categories; thus, $\mathsf{Set}$ will denote the category of sets. Thus
  \begin{itemize}
  \item $\operatorname{Obj}(\mathsf{Set}) = $ the class of all sets;
  \item for $A$, $B$ in $\operatorname{Obj}(\mathsf{Set})$ (that is, for $A$,
    $B$ sets) $\operatorname{Hom}_{\mathsf{Set}}(A, B) = B^A$.
  \end{itemize}
  Note that the presence of the operations recalled in \S{}\S{}1.3--1.5 is not
  part of the definition of category: these operations highlight interesting
  features of $\mathsf{Set}$, which may or may not be shared by other
  categories. We will soon come back to some of these operations and understand
  more precisely what they say about $\mathsf{Set}$.
\end{exam}
\begin{exam}{}{}
  Here is a completely different example. Suppose $S$ is a set and $\sim$ is a
  relation on $S$ satisfying the reflexive and transitive properties. Then we can
  encode this data into a category:
  \begin{itemize}
  \item \emph{objects}: the elements of $S$;
  \item \emph{morphisms}: if $a$, $b$ are objects (that is, if $a, b \in S$),
    then let $\operatorname{Hom}(a, b)$ be the set consisting of the element
    $(a, b) \in S \times S$ if $a \sim b$, and let $\operatorname{Hom}(a, b) =
    \emptyset$ otherwise.
  \end{itemize}
  Note that (unlike in $\mathsf{Set}$) there are very few morphisms: at most one
  for any pair of objects, and no morphisms at all between `unrelated' objects.

  We have to define `composition of morphisms' and verify that the conditions
  specified in \S{}3.1 are satisfied. First of all, do we have `identities'? If
  $a$ is an object (that is, if $a \in S$), we need to find an element
  $1_a \in \operatorname{Hom}(a, a)$. This is precisely why we are assuming that
  $\sim$ is reflexive: this tells us that $\forall a$, $a \sim a$; that is,
  $\operatorname{Hom}(a, a)$ consists of the single element $(a, a)$. So we have
  no choice: we must let
  \[ 1_a = (a, a) \in \operatorname{Hom}(a, a). \]
  As for composition, let $a$, $b$, $c$ be objects (that is, elements of $S$)
  and $f \in \operatorname{Hom}(a, b)$, $g \in \operatorname{Hom}(b, c)$; we
  have to define a corresponding morphism $gf \in \operatorname{Hom}(a, c)$.
  Now, $f \in \operatorname{Hom}(a, b)$ tells us that $\operatorname{Hom}(a, b)$
  is nonempty, and according to the definition of morphisms in this category that
  means $a \sim b$, and $f$ is in fact the element $(a, b)$ of $S \times S$.
  Similarly, $g \in \operatorname{Hom}(b, c)$ tells us $b \sim c$ and $g =
  (b, c)$. Now
  \[ a \sim b \text{ and } b \sim c \implies a \sim c \]
  since we are assuming that $\sim$ is transitive. This tells us that
  $\operatorname{Hom}(a, c)$ consists of the single element $(a, c)$. Thus we
  again have no choice: we must let
  \[ gf := (a, c) \in \operatorname{Hom}(a, c). \]
  Is this operation associative? If $f \in \operatorname{Hom}(a, b)$, $g \in
  \operatorname{Hom}(b, c)$, and $h \in \operatorname{Hom}(c, d)$, then
  necessarily $f = (a,b)$, $g = (b,c)$, $h = (c,d)$, and
  $gf = (a,c)$, $hg = (b,d)$, and hence $h(gf) = (a,d) = (hg)f$,
  proving associativity. The reader will have no difficulties checking that $1_a$
  is an identity with respect to this composition, as needed (Exercise~3.3).

  The most trivial instance of this construction is the category obtained from a
  set $S$ taken with the equivalence relation `$=$'; that is, the only morphisms
  are the identity morphisms. These categories are called \emph{discrete}.

  As another example, consider the category corresponding to endowing $\Z$ with
  the relation $\leq$. For example,
  \[
    \begin{tikzcd}
      2 \arrow[r] & 3 \arrow[r] \arrow[d, "1_3"'] & 5 \arrow[r] & 7 \\
      & 3 \arrow[urr] \arrow[r] & 4 \arrow[ur] &
    \end{tikzcd}
  \]
  is a (randomly chosen) commutative diagram in this category. It would still be
  a (commutative) diagram in this category if we reversed the vertical arrow
  $3 \to 3$ or if we added an arrow from $3$ to $4$, while we are not allowed to
  draw an arrow from $4$ to $3$, since $4 \not\leq 3$.

  These categories are very special---for example, every diagram drawn in them is
  necessarily commutative, and this is very far from being the case in, e.g.,
  $\mathsf{Set}$. Also note that these categories are all small.
\end{exam}
\begin{exam}{}{}
  Here is another\footnote{Actually, this is again an instance of the categories
    considered in Example~3.3. Do you see why? (Exercise~3.5.)} example of a
  small category. Let $S$ again be a set. Define a category $\hat{\mathsf{S}}$
  by setting
  \begin{itemize}
  \item $\operatorname{Obj}(\hat{\mathsf{S}}) = \mathscr{P}(S)$, the power set
    of $S$ (cf.\ \S{}1.2 and Exercise~2.11);
  \item for $A$, $B$ objects of $\hat{\mathsf{S}}$ (that is, $A \subseteq S$
    and $B \subseteq S$) let $\operatorname{Hom}_{\hat{\mathsf{S}}}(A, B)$ be
    the pair $(A, B)$ if $A \subseteq B$, and let
    $\operatorname{Hom}_{\hat{\mathsf{S}}}(A, B) = \emptyset$ otherwise.
  \end{itemize}
  The identity $1_A$ consists of the pair $(A, A)$ (which is one, and in fact
  the only one, morphism from $A$ to $A$, since $A \subseteq A$). Composition is
  obtained by stringing inclusions: if there are morphisms $A \to B$, $B \to C$
  in $\hat{\mathsf{S}}$, then $A \subseteq B$ and $B \subseteq C$; hence $A
  \subseteq C$ and there is a morphism $A \to C$. Checking the axioms specified
  in \S{}3.1 should be routine (make sure this is the case!).

  Examples in this style (but employing more sophisticated structures, such as
  the family of open subsets of a topological space) are hugely important in
  well-established fields such as algebraic geometry.
\end{exam}
\begin{exam}{}{}
  The next example is very abstract, but thinking about it will make you rather
  comfortable with everything we have seen so far; and it is a very common
  construction, variations of which will abound in the course.

  Let $\mathsf{C}$ be a category, and let $A$ be an object of $\mathsf{C}$. We
  are going to define a category $\mathsf{C}_A$ whose objects are certain
  morphisms in $\mathsf{C}$ and whose morphisms are certain diagrams of
  $\mathsf{C}$ (surprise!).
  \begin{itemize}
  \item $\operatorname{Obj}(\mathsf{C}_A) = $ all morphisms from any object of
    $\mathsf{C}$ to $A$; thus, an object of $\mathsf{C}_A$ is a morphism $f \in
    \operatorname{Hom}_{\mathsf{C}}(Z, A)$ for some object $Z$ of $\mathsf{C}$.
    Pictorially, an object of $\mathsf{C}_A$ is an arrow $Z \xrightarrow{f} A$
    in $\mathsf{C}$; these are often drawn `top-down', as in
    \[
      \begin{tikzcd} Z \arrow[d, "f"] \\ A \end{tikzcd}
    \]
  \end{itemize}
  What are morphisms in $\mathsf{C}_A$ going to be? There really is only one
  sensible way to assign morphisms to a category with objects as above. The brave
  reader will want to stop reading here and continue only after having come up
  with the definition independently. There will be many similar examples lurking
  behind constructions we will encounter in this book, and ideally speaking, they
  should appear completely natural when the time comes. A bit of effort devoted
  now to understanding this prototype situation will have ample reward in the
  future. Spoiler follows, so put these notes away now and jot down the
  definition of morphism in $\mathsf{C}_A$.

  Welcome back.
  \begin{itemize}
  \item Let $f_1$, $f_2$ be objects of $\mathsf{C}_A$, that is, two arrows
    $Z_1 \xrightarrow{f_1} A$ and $Z_2 \xrightarrow{f_2} A$ in $\mathsf{C}$.
    Morphisms $f_1 \to f_2$ are defined to be commutative diagrams
    \[
      \begin{tikzcd}
        Z_1 \arrow[rr, "\sigma"] \arrow[dr, "f_1"'] & & Z_2 \arrow[dl, "f_2"] \\
        & A &
      \end{tikzcd}
    \]
    in the `ambient' category $\mathsf{C}$.
  \end{itemize}
  That is, morphisms $f_1 \to f_2$ correspond precisely to those morphisms
  $\sigma : Z_1 \to Z_2$ in $\mathsf{C}$ such that $f_1 = f_2 \sigma$.

  Once you understand what morphisms have to be, checking that they satisfy the
  axioms spelled out in \S{}3.1 is straightforward. The identities are inherited
  from the identities in $\mathsf{C}$: for $f : Z \to A$ in $\mathsf{C}_A$, the
  identity $1_f$ corresponds to the diagram
  \[
    \begin{tikzcd}
      Z \arrow[rr, "1_Z"] \arrow[dr, "f"'] & & Z \arrow[dl, "f"] \\
      & A &
    \end{tikzcd}
  \]
  which commutes by virtue of the fact that $\mathsf{C}$ is a category.
  Composition is also a subproduct of composition in $\mathsf{C}$. Two morphisms
  $f_1 \to f_2 \to f_3$ in $\mathsf{C}_A$ correspond to putting two commutative
  diagrams side-by-side:
  \[
    \begin{tikzcd}
      Z_1 \arrow[r, "\sigma"] \arrow[dr, "f_1"'] &
      Z_2 \arrow[r, "\tau"] \arrow[d, "f_2"' description] &
      Z_3 \arrow[dl, "f_3"] \\
      & A &
    \end{tikzcd}
  \]
  and then it follows (again because $\mathsf{C}$ is a category!) that the
  diagram obtained by removing the central arrow, i.e.,
  \[
    \begin{tikzcd}
      Z_1 \arrow[rr, "\tau\sigma"] \arrow[dr, "f_1"'] & & Z_3 \arrow[dl, "f_3"] \\
      & A &
    \end{tikzcd}
  \]
  also commutes. Check all this(!), and verify that composition in $\mathsf{C}_A$
  is associative (again, this follows immediately from the fact that composition
  is associative in $\mathsf{C}$).

  Categories constructed in this fashion are called \emph{slice categories} in
  the literature; they are particular cases of \emph{comma categories}.
\end{exam}
\begin{exam}{}{}
  For the sake of concreteness, let's apply the construction given in Example~3.5
  to the category constructed in Example~3.3, say for $S = \Z$ and $\sim$ the
  relation $\leq$. Call $\mathsf{C}$ this category, and choose an object $A$ of
  $\mathsf{C}$---that is, an integer, for example, $A = 3$. Then the objects of
  $\mathsf{C}_A$ are morphisms in $\mathsf{C}$ with target $3$, that is, pairs
  $(n, 3) \in \Z \times \Z$ with $n \leq 3$. There is a morphism $(m, 3) \to
  (n, 3)$ if and only if $m \leq n$. In this case $\mathsf{C}_A$ may be
  harmlessly identified with the `subcategory' of integers $\leq 3$, with `the
  same' morphisms as in $\mathsf{C}$.
\end{exam}

\begin{exam}{}{}
  An entirely similar example to the one explored in Example~3.5 may be obtained
  by considering morphisms in a category $\mathsf{C}$ \emph{from} a fixed object
  $A$ to all objects in $\mathsf{C}$, again with morphisms defined by suitable
  commutative diagrams. This leads to \emph{coslice categories}. The reader
  should provide details of this construction (Exercise~3.7).
\end{exam}

\begin{exam}{}{}
  As a `concrete' instance of a category as in Example~3.7, let $\mathsf{C} =
  \mathsf{Set}$ and $A = $ a fixed singleton $\{*\}$. Call the resulting category
  $\mathsf{Set}^*$.

  An object in $\mathsf{Set}^*$ is then a morphism $f : \{*\} \to S$ in
  $\mathsf{Set}$, where $S$ is any set. The information of an object in
  $\mathsf{Set}^*$ consists therefore of the choice of a nonempty set $S$ and of
  an element $s \in S$---that is, the element $f(*)$: this element determines,
  and is determined by, $f$.

  Thus, we may denote objects of $\mathsf{Set}^*$ as pairs $(S, s)$, where $S$
  is any set and $s \in S$ is any element of $S$. A morphism between two such
  objects, $(S, s) \to (T, t)$, corresponds then (check this!) to a set-function
  $\sigma : S \to T$ such that $\sigma(s) = t$.

  Objects of $\mathsf{Set}^*$ are called \emph{pointed sets}. Many of the
  structures we will study in this book will be pointed sets. For example (as we
  will see) a `group' is a set $G$ with, among other requirements, a
  distinguished element $e_G$ (its `identity'); `group homomorphisms' will be
  functions which, among other properties, send identities to identities; thus,
  they are morphisms of pointed sets in the sense considered above.
\end{exam}

\begin{exam}{}{}
  It is useful to contemplate a few more `abstract' examples in the style of
  Examples~3.5 and~3.7. These will be essential ingredients in the promised
  revisitation of some of the operations mentioned in \S{}1.3. Their definition
  will appear disappointedly simple-minded to the reader who has mastered
  Examples~3.5 and~3.7.

  This time we start from a given category $\mathsf{C}$ and \emph{two} objects
  $A$, $B$ of $\mathsf{C}$. We can define a new category $\mathsf{C}_{A,B}$ by
  essentially the same procedure that we used in order to define $\mathsf{C}_A$:
  \begin{itemize}
  \item $\operatorname{Obj}(\mathsf{C}_{A,B})$ = diagrams
    % TODO: commutative diagram with Z mapping to both A and B (two-legged span)
    $Z \xrightarrow{f} A$ and $Z \xrightarrow{g} B$ in $\mathsf{C}$; and
  \item morphisms from $(f_1 : Z_1 \to A,\, g_1 : Z_1 \to B)$ to
    $(f_2 : Z_2 \to A,\, g_2 : Z_2 \to B)$ are those morphisms $\sigma : Z_1
    \to Z_2$ in $\mathsf{C}$ making both diagrams commute simultaneously:
    $f_1 = f_2 \sigma$ and $g_1 = g_2 \sigma$.
  \end{itemize}
  I will leave to the reader the task of formalizing this rough description. This
  example is really nothing more than a mixture of $\mathsf{C}_A$ and
  $\mathsf{C}_B$, where the two structures interact because of the stringent
  requirement that the same $\sigma$ must make both sides of the diagram commute
  simultaneously.

  Flipping most of the arrows gives an analogous variation of Example~3.7,
  producing a category which we may\footnote{There does not seem to be an
    established notation for these commonplace categories.} denote
  $\mathsf{C}^{A,B}$; details are left to the reader.
\end{exam}

\begin{exam}{}{}
  As a final variation on these examples, we conclude by considering the
  \emph{fibered} version of $\mathsf{C}_{A,B}$ (and $\mathsf{C}^{A,B}$). Take
  this as a test to see if you have really understood $\mathsf{C}_{A,B}$---experts
  would tell you that this looks fairly sophisticated for students just learning
  categories, so don't get disheartened if it does not flow too well at first
  (but pat yourself on the shoulder if it does!). Start with a given category
  $\mathsf{C}$, and this time choose two fixed morphisms $\alpha : A \to C$,
  $\beta : B \to C$ in $\mathsf{C}$, with the same target $C$. We can then
  consider a category $\mathsf{C}_{\alpha, \beta}$ as follows:
  \begin{itemize}
  \item $\operatorname{Obj}(\mathsf{C}_{\alpha,\beta})$ = commutative diagrams
    % TODO: diagram with Z mapping to A and B, composed with α and β to a common C
    in $\mathsf{C}$: pairs $(f : Z \to A,\, g : Z \to B)$ such that $\alpha f =
    \beta g$; and
  \item morphisms correspond to morphisms $\sigma : Z_1 \to Z_2$ making the
    evident diagram commute.
  \end{itemize}
  A solid understanding of Example~3.9 will make this example look just as tame;
  at this point the reader should have no difficulties formalizing it (that is,
  explaining how composition works, what identities are, etc.).

  Also left to the reader is the construction of the `mirror' example
  $\mathsf{C}^{\alpha,\beta}$, starting from two morphisms $\alpha : C \to A$,
  $\beta : C \to B$ with common source.
\end{exam}

\subsection*{Exercises}

\begin{exer}{}{}
  \exerused Let $\mathsf{C}$ be a category. Consider a structure $\mathsf{C}^{op}$ with
  \begin{itemize}
  \item $\operatorname{Obj}(\mathsf{C}^{op}) := \operatorname{Obj}(\mathsf{C})$;
  \item for $A$, $B$ objects of $\mathsf{C}^{op}$ (hence objects of $\mathsf{C}$),
    $\operatorname{Hom}_{\mathsf{C}^{op}}(A, B) := \operatorname{Hom}_{\mathsf{C}}(B, A)$.
  \end{itemize}
  Show how to make this into a category (that is, define composition of morphisms
  in $\mathsf{C}^{op}$ and verify the properties listed in \S{}3.1). Intuitively,
  the `opposite' category $\mathsf{C}^{op}$ is simply obtained by `reversing all
  the arrows' in $\mathsf{C}$. [5.1, \S{}VIII.1.1, \S{}IX.1.2, IX.1.10]
\end{exer}

\begin{exer}{}{}
  If $A$ is a finite set, how large is $\operatorname{End}_{\mathsf{Set}}(A)$?
\end{exer}

\begin{exer}{}{}
  \exerused Formulate precisely what it means to say that $1_a$ is an identity
  with respect to composition in Example~3.3, and prove this assertion. [\S{}3.2]
\end{exer}

\begin{exer}{}{}
  Can we define a category in the style of Example~3.3 using the relation $<$ on
  the set $\Z$?
\end{exer}

\begin{exer}{}{}
  \exerused Explain in what sense Example~3.4 is an instance of the categories
  considered in Example~3.3. [\S{}3.2]
\end{exer}

\begin{exer}{}{}
  \exerused (Assuming some familiarity with linear algebra.) Define a category
  $\mathsf{V}$ by taking $\operatorname{Obj}(\mathsf{V}) = \N$ and letting
  $\operatorname{Hom}_{\mathsf{V}}(n, m) = $ the set of $m \times n$ matrices
  with real entries, for all $n, m \in \N$. (I will leave the reader the task of
  making sense of a matrix with 0 rows or columns.) Use product of matrices to
  define composition. Does this category `feel' familiar? [\S{}VI.2.1,
  \S{}VIII.1.3]
\end{exer}

\begin{exer}{}{}
  \exerused Define carefully objects and morphisms in Example~3.7, and draw the
  diagram corresponding to composition. [\S{}3.2]
\end{exer}

\begin{exer}{}{}
  \exerused A \emph{subcategory} $\mathsf{C}'$ of a category $\mathsf{C}$
  consists of a collection of objects of $\mathsf{C}$ with sets of morphisms
  $\operatorname{Hom}_{\mathsf{C}'}(A, B) \subseteq \operatorname{Hom}_{\mathsf{C}}(A, B)$
  for all objects $A$, $B$ in $\operatorname{Obj}(\mathsf{C}')$, such that
  identities and compositions in $\mathsf{C}$ make $\mathsf{C}'$ into a
  category. A subcategory $\mathsf{C}'$ is \emph{full} if
  $\operatorname{Hom}_{\mathsf{C}'}(A, B) = \operatorname{Hom}_{\mathsf{C}}(A, B)$
  for all $A$, $B$ in $\operatorname{Obj}(\mathsf{C}')$. Construct a category of
  infinite sets and explain how it may be viewed as a full subcategory of
  $\mathsf{Set}$. [4.4, \S{}VI.1.1, \S{}VIII.1.3]
\end{exer}

\begin{exer}{}{}
  \exerused An alternative to the notion of multiset introduced in \S{}2.2 is
  obtained by considering sets endowed with equivalence relations; equivalent
  elements are taken to be multiple instances of elements `of the same kind'.
  Define a notion of morphism between such enhanced sets, obtaining a category
  $\mathsf{MSet}$ containing (a `copy' of) $\mathsf{Set}$ as a full subcategory.
  (There may be more than one reasonable way to do this! This is intentionally an
  open-ended exercise.) Which objects in $\mathsf{MSet}$ determine ordinary
  multisets as defined in \S{}2.2 and how? Spell out what a morphism of multisets
  would be from this point of view. (There are several natural notions of
  morphisms of multisets. Try to define morphisms in $\mathsf{MSet}$ so that the
  notion you obtain for ordinary multisets captures your intuitive understanding
  of these objects.) [\S{}2.2, \S{}3.2, 4.5]
\end{exer}

\begin{exer}{}{}
  Since the objects of a category $\mathsf{C}$ are not (necessarily interpreted
  as) sets, it is not clear how to make sense of a notion of `subobject' in
  general. In some situations it does make sense to talk about subobjects, and
  the subobjects of any given object $A$ in $\mathsf{C}$ are in one-to-one
  correspondence with the morphisms $A \to \Omega$ for a fixed, special object
  $\Omega$ of $\mathsf{C}$, called a \emph{subobject classifier}. Show that
  $\mathsf{Set}$ has a subobject classifier.
\end{exer}

\begin{exer}{}{}
  \exerused Draw the relevant diagrams and define composition and identities for
  the category $\mathsf{C}^{A,B}$ mentioned in Example~3.9. Do the same for the
  category $\mathsf{C}^{\alpha,\beta}$ mentioned in Example~3.10.
  [\S{}5.5, 5.12]
\end{exer}

\section{Morphisms}
\label{sec:prelims:morphisms}

Just as in $\mathsf{Set}$ we highlight certain types of functions (injective,
surjective, bijective), it is useful to try to do the same for morphisms in an
arbitrary category. The reader should note that defining qualities of morphisms
by their actions on `elements' is not an option in the general setting, because
objects of an arbitrary category do not (in general) have `elements'.

This is why we spent some time analyzing injectivity, etc., from different
viewpoints in \S{}\S{}2.4--2.6. It turns out that the other viewpoints on these
notions do transfer nicely into the categorical setting.

\subsection{Isomorphisms}
\label{subsec:prelims:isomorphisms}
Let $\mathsf{C}$ be a category.

\begin{defn}{Isomorphism}{defn:prelims:isomorphism}
A morphism $f \in \operatorname{Hom}_{\mathsf{C}}(A,B)$ is an
\textbf{isomorphism} if it has a (two-sided) inverse under composition: that
is, if $\exists\, g \in \operatorname{Hom}_{\mathsf{C}}(B,A)$ such that
\[
  g \circ f = 1_A, \qquad f \circ g = 1_B.
\]
\end{defn}

Recall that in \S{}2.5 the inverse of a bijection of sets $f$ was defined
`elementwise'; in particular, there was no ambiguity in its definition, and we
introduced the notation $f^{-1}$ for this function. By contrast, the `inverse'
$g$ produced in \Cref{defn:prelims:isomorphism} does not appear to have this
uniqueness explicitly built into its definition. Luckily, its defining property
does guarantee its uniqueness, but this requires a verification:

\begin{prop}{The inverse of an isomorphism is unique.}{prop:prelims:iso-unique-inverse}
\end{prop}
\begin{proof}
We have to verify that if both $g_1$ and $g_2 \colon B \to A$ act as inverses
of a given isomorphism $f \colon A \to B$, then $g_1 = g_2$. The standard
trick for this kind of verification is to compose $f$ on the left by one of the
morphisms, and on the right by the other; then apply associativity. The whole
argument can be compressed into one line:\footnote{Associativity of composition
implies that parentheses may be shuffled at will in longer expressions, as done
here (cf.\ \Cref{exer:4.1}).}
\[
  g_1 = g_1 \circ 1_B = g_1 \circ (f \circ g_2) = (g_1 \circ f) \circ g_2
      = 1_A \circ g_2 = g_2,
\]
as needed.
\end{proof}

Note that the argument really proves that if $f$ is a morphism with a
left-inverse $g_1$ and a right-inverse $g_2$, then necessarily $f$ is an
isomorphism, $g_1 = g_2$, and this morphism is the (unique) inverse of $f$.
Look back at Corollary~2.2. Since the inverse of $f$ is uniquely determined by
$f$, there is no ambiguity in denoting it by $f^{-1}$.

\begin{prop}{}{prop:prelims:iso-props}
With notation as above:
\begin{itemize}
  \item Each identity $1_A$ is an isomorphism and is its own inverse.
  \item If $f$ is an isomorphism, then $f^{-1}$ is an isomorphism and
        $(f^{-1})^{-1} = f$.
  \item If $f \in \operatorname{Hom}_{\mathsf{C}}(A,B)$ and
        $g \in \operatorname{Hom}_{\mathsf{C}}(B,C)$ are isomorphisms, then
        $g \circ f$ is an isomorphism and $(g \circ f)^{-1} = f^{-1} \circ g^{-1}$.
\end{itemize}
\end{prop}
\begin{proof}
These all `prove themselves'. For example, it is immediate to verify that
$f^{-1} \circ g^{-1}$ is a left-inverse of $g \circ f$: indeed,
\[
  (f^{-1} \circ g^{-1}) \circ (g \circ f)
    = f^{-1} \circ \bigl((g^{-1} \circ g) \circ f\bigr)
    = f^{-1} \circ (1_B \circ f)
    = f^{-1} \circ f = 1_A.
\]
The verification that $f^{-1} \circ g^{-1}$ is also a right-inverse of
$g \circ f$ is analogous.
\end{proof}

Note that taking the inverse reverses the order of composition:
$(g \circ f)^{-1} = f^{-1} \circ g^{-1}$.

Two objects $A$, $B$ of a category are \textbf{isomorphic} if there is an
isomorphism $f \colon A \to B$. An immediate corollary of
\Cref{prop:prelims:iso-props} is that `isomorphic' is an equivalence
relation.\footnote{The reader should have checked this in \Cref{exer:2.4}, for
$\mathsf{Set}$; the same proof will work in any category.} If two objects $A$,
$B$ are isomorphic, one writes $A \cong B$.

\begin{exam}{}{}
Of course, the isomorphisms in the category $\mathsf{Set}$ are precisely the
bijections; this was observed at the beginning of \S{}2.5.
\end{exam}

\begin{exam}{}{}
As noted in \Cref{prop:prelims:iso-props}, identities are isomorphisms. They
may be the only isomorphisms in a category: for example, this is the case in
the category $\mathsf{C}$ obtained from the relation $\leq$ on $\mathbb{Z}$,
as in \Cref{exam:prelims:category-leq}. Indeed, for $a$, $b$ objects of
$\mathsf{C}$ (that is, $a, b \in \mathbb{Z}$), there is a morphism
$f \colon a \to b$ and a morphism $g \colon b \to a$ only if $a \leq b$ and
$b \leq a$, that is, if $a = b$. So an isomorphism in $\mathsf{C}$ necessarily
acts from an object $a$ to itself; but in $\mathsf{C}$ there is only one such
morphism, that is, $1_a$.
\end{exam}

\begin{exam}{}{}
On the other hand, there are categories in which every morphism is an
isomorphism; such categories are called \textbf{groupoids}. The reader
`already knows' many examples of groupoids; cf.\ \Cref{exer:4.2}.
\end{exam}

An \textbf{automorphism} of an object $A$ of a category $\mathsf{C}$ is an
isomorphism from $A$ to itself. The set of automorphisms of $A$ is denoted
$\operatorname{Aut}_{\mathsf{C}}(A)$; it is a subset of
$\operatorname{End}_{\mathsf{C}}(A)$. By \Cref{prop:prelims:iso-props},
composition confers on $\operatorname{Aut}_{\mathsf{C}}(A)$ a remarkable
structure:
\begin{itemize}
  \item the composition of two elements $f, g \in \operatorname{Aut}_{\mathsf{C}}(A)$
        is an element $g \circ f \in \operatorname{Aut}_{\mathsf{C}}(A)$;
  \item composition is associative;
  \item $\operatorname{Aut}_{\mathsf{C}}(A)$ contains the element $1_A$, which
        is an identity for composition (that is,
        $f \circ 1_A = 1_A \circ f = f$);
  \item every element $f \in \operatorname{Aut}_{\mathsf{C}}(A)$ has an inverse
        $f^{-1} \in \operatorname{Aut}_{\mathsf{C}}(A)$.
\end{itemize}
In other words, $\operatorname{Aut}_{\mathsf{C}}(A)$ is a group, for all
objects $A$ of all categories $\mathsf{C}$. We will soon devote all our
attention to groups!

\subsection{Monomorphisms and epimorphisms}
\label{subsec:prelims:morphisms:monomorphisms-and-epimorphisms}
As pointed out above, we do not have the option of defining for morphisms of an
arbitrary category a notion such as `injective' in the same way as we do for
set-functions in \S{}2.4: that definition requires a notion of `element', and
in general no such notion is available for objects of a category. But nothing
prevents us from defining monomorphisms as we did in \S{}2.6, in an arbitrary
category:

\begin{defn}{Monomorphism}{defn:prelims:monomorphism}
Let $\mathsf{C}$ be a category. A morphism $f \in \operatorname{Hom}_{\mathsf{C}}(A,B)$
is a \textbf{monomorphism} if the following holds: for all objects $Z$ of
$\mathsf{C}$ and all morphisms $\alpha', \alpha'' \in
\operatorname{Hom}_{\mathsf{C}}(Z,A)$,
\[
  f \circ \alpha' = f \circ \alpha'' \implies \alpha' = \alpha''.
\]
\end{defn}

Similarly, epimorphisms are defined as follows:

\begin{defn}{Epimorphism}{defn:prelims:epimorphism}
Let $\mathsf{C}$ be a category. A morphism $f \in \operatorname{Hom}_{\mathsf{C}}(A,B)$
is an \textbf{epimorphism} if the following holds: for all objects $Z$ of
$\mathsf{C}$ and all morphisms $\beta', \beta'' \in
\operatorname{Hom}_{\mathsf{C}}(B,Z)$,
\[
  \beta' \circ f = \beta'' \circ f \implies \beta' = \beta''.
\]
\end{defn}

\begin{exam}{}{}
As proven in Proposition~2.3, in the category $\mathsf{Set}$ the monomorphisms
are precisely the injective functions. The reader should have by now checked
that, likewise, in $\mathsf{Set}$ the epimorphisms are precisely the surjective
functions (cf.\ \Cref{exer:2.5}). Thus, while the definitions given in \S{}2.6
may have looked counterintuitive at first, they work as natural `categorical
counterparts' of the ordinary notions of injective/surjective functions.
\end{exam}

\begin{exam}{}{}
In the categories of \Cref{exam:prelims:category-leq}, every morphism is both
a monomorphism and an epimorphism. Indeed, recall that there is at most one
morphism between any two objects in these categories; hence the conditions
defining monomorphisms and epimorphisms are vacuous.
\end{exam}

Contemplating the previous example reveals a few unexpected twists in these
definitions, which defy our intuition as set-theorists. For instance, in
$\mathsf{Set}$, a function is an isomorphism if and only if it is both injective
and surjective, hence if and only if it is both a monomorphism and an
epimorphism. But in the category defined by $\leq$ on $\mathbb{Z}$, every
morphism is both a monomorphism and an epimorphism, while the only isomorphisms
are the identities (\Cref{exam:prelims:category-leq}). Thus this property is a
special feature of $\mathsf{Set}$, and we should not expect it to hold
automatically in every category; it will not hold in the category $\mathsf{Ring}$
of rings. It will hold in every abelian category (of which $\mathsf{Set}$ is
not an example!), but that is a story for a very distant future.

Similarly, in $\mathsf{Set}$ a function is an epimorphism, that is, surjective,
if and only if it has a right-inverse (Proposition~2.1); this may fail in
general, even in respectable categories such as the category $\mathsf{Grp}$ of
groups.

\subsection*{Exercises}

\begin{exer}{\exerused}{exer:4.1}
Composition is defined for two morphisms. If more than two morphisms are given,
e.g.,
\[
  A \xrightarrow{f} B \xrightarrow{g} C \xrightarrow{h} D \xrightarrow{i} E,
\]
then one may compose them in several ways, for example: $(i \circ h) \circ
(g \circ f)$, $(i \circ (h \circ g)) \circ f$, $i \circ ((h \circ g) \circ f)$,
etc., so that at every step one is only composing two morphisms. Prove that the
result of any such nested composition is independent of the placement of the
parentheses.

\begin{hint}
Use induction on $n$ to show that any such choice for $f_n \circ f_{n-1} \circ
\cdots \circ f_1$ equals $((\cdots((f_n \circ f_{n-1}) \circ f_{n-2}) \circ
\cdots) \circ f_1)$. Carefully working out the case $n = 5$ is helpful.
\end{hint}
[\S{}4.1, \S{}II.1.3]
\end{exer}

\begin{exer}{\exerused}{exer:4.2}
In \Cref{exam:prelims:category-leq} we have seen how to construct a category
from a set endowed with a relation, provided this latter is reflexive and
transitive. For what types of relations is the corresponding category a groupoid
(cf.\ \Cref{exam:prelims:groupoids})? [\S{}4.1]
\end{exer}

\begin{exer}{}{exer:4.3}
Let $A$, $B$ be objects of a category $\mathsf{C}$, and let
$f \in \operatorname{Hom}_{\mathsf{C}}(A,B)$ be a morphism.
\begin{itemize}
  \item Prove that if $f$ has a right-inverse, then $f$ is an epimorphism.
  \item Show that the converse does not hold, by giving an explicit example of a
        category and an epimorphism without a right-inverse.
\end{itemize}
\end{exer}

\begin{exer}{}{exer:4.4}
Prove that the composition of two monomorphisms is a monomorphism. Deduce that
one can define a subcategory $\mathsf{C}_{\mathrm{mono}}$ of a category
$\mathsf{C}$ by taking the same objects as in $\mathsf{C}$ and defining
$\operatorname{Hom}_{\mathsf{C}_{\mathrm{mono}}}(A,B)$ to be the subset of
$\operatorname{Hom}_{\mathsf{C}}(A,B)$ consisting of monomorphisms, for all
objects $A$, $B$. (Cf.\ \Cref{exer:3.8}; of course, in general
$\mathsf{C}_{\mathrm{mono}}$ is not full in $\mathsf{C}$.) Do the same for
epimorphisms. Can you define a subcategory $\mathsf{C}_{\mathrm{nonmono}}$ of
$\mathsf{C}$ by restricting to morphisms that are \emph{not} monomorphisms?
\end{exer}

\begin{exer}{}{exer:4.5}
Give a concrete description of monomorphisms and epimorphisms in the category
$\mathsf{MSet}$ you constructed in \Cref{exer:3.9}. (Your answer will depend
on the notion of morphism you defined in that exercise!)
\end{exer}

\section{Universal properties}
\label{sec:prelims:universal-properties}

The `abstract' examples in \S{}3 may have left the reader with the impression
that one can produce at will a large number of minute variations of the same
basic ideas, without really breaking any new ground. This may be fun in itself,
but \emph{why} do we really want to explore this territory?
Categories offer a rich unifying language, giving us a bird's eye view of many
constructions in algebra (and other fields). In this course, this will be most
apparent in the steady appearance of constructions satisfying suitable
\emph{universal properties}. For instance, we will see in a moment that products
and disjoint unions (as reviewed in \S{}1.3 and following) are characterized by
certain universal properties having to do with the categories $\mathsf{C}_{A,B}$
and $\mathsf{C}^{A,B}$ considered in Example~3.9.
Many of the concepts introduced in this course will have an explicit description
(such as the definition of product of sets given in \S{}1.4) and an accompanying
description in terms of a universal property (such as the one we will see in
\S{}5.4). The `explicit' description may be very useful in concrete computations
or arguments, but as a rule it is the universal property that clarifies the true
nature of the construction. In some cases (such as for the disjoint union) the
explicit description may turn out to depend on a seemingly arbitrary choice,
while the universal property will have no element of arbitrariness. In fact,
viewing the construction in terms of its corresponding universal property
clarifies why one can only expect it to be defined `up to isomorphism'.
Also, deeper relationships become apparent when the constructions are viewed in
terms of their universal properties. For example, we will see that products of
sets and disjoint unions of sets are really `mirror' constructions (in the sense
that reversing arrows transforms the universal property for one into that for the
other). This is not so clear (to this writer, anyway) from the explicit
descriptions in \S{}1.4.

\subsection{Initial and final objects}
\label{subsec:prelims:initial-and-final-objects}

\begin{defn}{Initial and final objects}{defn:prelims:initial-final}
Let $\mathsf{C}$ be a category. We say that an object $I$ of $\mathsf{C}$ is
\textbf{initial} in $\mathsf{C}$ if for every object $A$ of $\mathsf{C}$ there
exists exactly one morphism $I \to A$ in $\mathsf{C}$:
\[
  \forall A \in \operatorname{Obj}(\mathsf{C}) \colon \quad
  \operatorname{Hom}_{\mathsf{C}}(I,A) \text{ is a singleton.}
\]
We say that an object $F$ of $\mathsf{C}$ is \textbf{final} in $\mathsf{C}$ if
for every object $A$ of $\mathsf{C}$ there exists exactly one morphism
$A \to F$ in $\mathsf{C}$:
\[
  \forall A \in \operatorname{Obj}(\mathsf{C}) \colon \quad
  \operatorname{Hom}_{\mathsf{C}}(A,F) \text{ is a singleton.}
\]
\end{defn}

One may use \emph{terminal} to denote either possibility, but in general I would
advise the reader to be explicit about which `end' of $\mathsf{C}$ one is
considering.

A category need not have initial or final objects, as the following example
shows.
\begin{exam}{}{}
The category obtained by endowing $\mathbb{Z}$ with the relation $\leq$ (see
Example~3.3) has no initial or final object. Indeed, an initial object in this
category would be an integer $i$ such that $i \leq a$ for all integers $a$;
there is no such integer. Similarly, a final object would be an integer $f$
larger than every integer, and there is no such thing.

By contrast, the category considered in Example~3.6 \emph{does} have a final
object, namely the pair $(3,3)$; it still has no initial object.
\end{exam}
Also, initial and final objects, when they exist, may or may not be unique:

\begin{exam}{}{}
In $\mathsf{Set}$, the empty set $\emptyset$ is initial (the `empty graph'
defines the unique function from $\emptyset$ to any given object!), and clearly
it is the unique set that fits this requirement (\Cref{exer:5.2}).

$\mathsf{Set}$ also has final objects: for every set $A$, there is a unique
function from $A$ to a singleton $\{p\}$ (that is, the `constant' function).
Every singleton is final in $\mathsf{Set}$; thus, final objects are not unique
in this category.
\end{exam}
However, I claim that if initial/final objects exist, then they are unique
\emph{up to a unique isomorphism}. I will invoke this fact frequently, so here
is its official statement and its (immediate) proof:

\begin{prop}{}{prop:prelims:initial-final-unique}
Let $\mathsf{C}$ be a category.
\begin{itemize}
  \item If $I_1$, $I_2$ are both initial objects in $\mathsf{C}$, then
        $I_1 \cong I_2$.
  \item If $F_1$, $F_2$ are both final objects in $\mathsf{C}$, then
        $F_1 \cong F_2$.
\end{itemize}
Further, these isomorphisms are uniquely determined.
\end{prop}
\begin{proof}
Recall that (by definition of category!) for every object $A$ of $\mathsf{C}$
there is at least one element in $\operatorname{Hom}_{\mathsf{C}}(A,A)$, namely
the identity $1_A$. If $I$ is initial, then there is a unique morphism
$I \to I$, which therefore must be the identity $1_I$.

Now assume $I_1$ and $I_2$ are both initial in $\mathsf{C}$. Since $I_1$ is
initial, there is a unique morphism $f \colon I_1 \to I_2$ in $\mathsf{C}$; we
have to show that $f$ is an isomorphism. Since $I_2$ is initial, there is a
unique morphism $g \colon I_2 \to I_1$ in $\mathsf{C}$. Consider
$g \circ f \colon I_1 \to I_1$; as observed, necessarily
\[
  g \circ f = 1_{I_1}
\]
since $I_1$ is initial. By the same token $f \circ g = 1_{I_2}$ since $I_2$ is
initial. This proves that $f \colon I_1 \to I_2$ is an isomorphism, as needed.
The proof for final objects is entirely analogous (\Cref{exer:5.3}).
\end{proof}

\Cref{prop:prelims:initial-final-unique} ``explains'' why, while not unique, the
final objects in $\mathsf{Set}$ are all isomorphic: no singleton is more
`special' than any other singleton; this is the typical situation. There may be
psychological reasons why one initial or final object looks more compelling than
others (for example, the singleton $\{\emptyset\} = 2^\emptyset$ may look to
some like the most `natural' choice among all singletons), but this plays no
role in how these objects sit in their category.

\subsection{Universal properties}
\label{subsec:prelims:universal-properties}
The most natural context in which to introduce \emph{universal properties}
requires a good familiarity with the language of functors, which we will only
introduce at a later stage (cf.\ \S{}VIII.1.1). For the purpose of the examples
we will run across in (most of) this book, the following `working definition'
should suffice.
We say that a construction \emph{satisfies a universal property} (or `is the
solution to a universal problem') when it may be viewed as a terminal object of
a category. The category depends on the context and is usually explained `in
words' (and often without even mentioning the word \emph{category}).
In particularly simple cases this may take the form of a statement such as
\emph{$\emptyset$ is universal with respect to the property of mapping to sets};
this is synonymous with the assertion that $\emptyset$ is initial in the
category $\mathsf{Set}$.
More often, the situation is more complex. Since being initial/final amounts to
the existence and uniqueness of certain morphisms, the `explanation' of a
universal property may follow the pattern, ``object $X$ is universal with
respect to the following property: for any $Y$ such that\ldots, there exists a
unique morphism $Y \to X$ such that\ldots''
The not-so-naive reader will recognize that this explanation hides the
definition of an accessory category and the statement that $X$ is terminal
(probably final in this case) in this new category. It is useful to learn how
to translate such wordy explanations into what they really mean. Also, the
reader should keep in mind that it is not uncommon to sweep under the rug part
of the essential information about the solution to a universal problem (usually
some key morphism): this information is presumably implicit in any given set-up.
This will be apparent from the examples that follow.

\subsection{Quotients}
\label{subsec:prelims:quotients}
Let $\sim$ be an equivalence relation defined on a set $A$. Let's parse the
assertion:
\begin{quote}
``The quotient $A/{\sim}$ is universal with respect to the property of mapping
$A$ to a set in such a way that equivalent elements have the same image.''
\end{quote}
What can this possibly mean, and is it true?
The assertion is talking about functions
\[
  A \xrightarrow{\varphi} Z
\]
with $Z$ any set, satisfying the property
\[
  a' \sim a'' \implies \varphi(a') = \varphi(a'').
\]
These morphisms are objects of a category (very similar to the category defined
in Example~3.7); for convenience, let's denote such an object by $(\varphi, Z)$.
The only reasonable way to define morphisms $(\varphi_1, Z_1) \to (\varphi_2, Z_2)$
is as commutative diagrams
\[
\begin{tikzcd}
  Z_1 \arrow[rr, "\sigma"] & & Z_2 \\
  & A \arrow[ul, "\varphi_1"] \arrow[ur, "\varphi_2"'] &
\end{tikzcd}
\]
This is the same definition considered in Example~3.7.

Does this category have initial objects?
\begin{prop}{}{prop:prelims:quotient-universal}
Denoting by $\pi$ the `canonical projection' defined in Example~2.6, the pair
$(\pi, A/{\sim})$ is an initial object of this category.
\end{prop}

This is what our writer meant by the mysterious assertion copied above. Once
this is understood, it is very easy to prove that the assertion is indeed
correct.

\begin{proof}
Consider any $(\varphi, Z)$ as above. We have to prove that there exists a
unique morphism $(\pi, A/{\sim}) \to (\varphi, Z)$, that is, a unique
commutative diagram
\[
\begin{tikzcd}
  A/{\sim} \arrow[rr, "\overline{\varphi}"] & & Z \\
  & A \arrow[ul, "\pi"] \arrow[ur, "\varphi"'] &
\end{tikzcd}
\]
that is, a unique function $\overline{\varphi}$ making this diagram commute.

Let $[a]_\sim$ be an arbitrary element of $A/{\sim}$. If the diagram is indeed
going to commute, then necessarily
\[
  \overline{\varphi}([a]_\sim) = \varphi(a);
\]
this tells us that $\overline{\varphi}$ is indeed unique, if it exists at
all---that is, if this prescription does define a function $A/{\sim} \to Z$.

Hence, all we have to check is that $\overline{\varphi}$ is well-defined, that
is, that if $[a_1]_\sim = [a_2]_\sim$, then $\varphi(a_1) = \varphi(a_2)$;
and indeed
\[
  [a_1]_\sim = [a_2]_\sim \implies a_1 \sim a_2 \implies \varphi(a_1) = \varphi(a_2).
\]
This is precisely the condition that morphisms in our category satisfy.
\end{proof}
Note the several levels of sloppiness in the assertion considered above: it
does not tell us very explicitly what category to consider; it does not tell us
that we should especially pay attention to initial objects in this category.
Worst of all, the solution to the universal problem is not really $A/{\sim}$,
but rather \emph{the morphism} $\pi \colon A \to A/{\sim}$.

The reader should practice the skill of translating loose assertions such as the
one given above into precise statements; it is not at all uncommon to run into
examples at the same level of `abuse of language' as this one.

The reason why we get away with writing such assertions is that the context
really allows the experienced reader to parse them effectively, and they are
substantially more concise than their spelled-out version. After all, there is
in general no conceivable choice for a morphism $A \to A/{\sim}$ other than the
canonical projection; hence, neglecting to mention it is forgivable. Also, the
final object in the category considered above is supremely uninteresting (what
is it? cf.\ \Cref{exer:5.5}), so surely we must have meant the initial one.
What do we learn by viewing quotients in terms of their universal property? For
example, suppose $\sim$ is the equivalence relation defined starting from a
function $f \colon A \to B$, as in \S{}2.8. Then the reader will realize easily
that $\image f$ also satisfies the universal property given above for $A/{\sim}$;
therefore (by \Cref{prop:prelims:initial-final-unique}) $\image f$ and $A/{\sim}$
must be isomorphic. This is precisely the content of Theorem~2.7; thus, the
universal property sheds some light on the `canonical decomposition' studied in
\S{}2.8.

\subsection{Products}
\label{subsec:prelims:products}
It is also a very good exercise to stare at a familiar construction and try to
see the universal property which may be behind it. I will now encourage you,
dear reader, to contemplate the notion of the product of two sets given in
\S{}1.4 and to see if the universal property it satisfies jumps out at you.
Spoiler follows, so this is a good time to stop reading these notes and to try
on your own.
Here is the universal property. Let $A$, $B$ be sets, and consider the product
$A \times B$, with the two natural projections:
\[
\begin{tikzcd}[column sep=small]
  & A \times B \arrow[dl, "\pi_A"'] \arrow[dr, "\pi_B"] & \\
  A & & B
\end{tikzcd}
\]
(see Example~2.4). Then for every set $Z$ and morphisms
\[
\begin{tikzcd}[column sep=small]
  & Z \arrow[dl, "f_A"'] \arrow[dr, "f_B"] & \\
  A & & B
\end{tikzcd}
\]
there exists a unique morphism $\sigma \colon Z \to A \times B$ such that the
diagram
\[
\begin{tikzcd}
  & Z \arrow[d, "\sigma"] \arrow[ddl, "f_A"'] \arrow[ddr, "f_B"] & \\
  & A \times B \arrow[dl, "\pi_A"] \arrow[dr, "\pi_B"'] & \\
  A & & B
\end{tikzcd}
\]
commutes. In this situation, $\sigma$ is usually denoted $f_A \times f_B$.

\begin{proof}
Define $\forall z \in Z$:
\[
  \sigma(z) = (f_A(z), f_B(z)).
\]
This function\footnote{Note that there is no `well-definedness' issue this
time.} manifestly makes the diagram commute: $\forall z \in Z$,
\[
  \pi_A \sigma(z) = \pi_A(f_A(z), f_B(z)) = f_A(z),
\]
showing that $\pi_A \circ \sigma = f_A$ and similarly $\pi_B \circ \sigma =
f_B$. Further, the definition is forced by the commutativity of the diagram;
so $\sigma$ is unique, as claimed.
\end{proof}
% In other words, products of sets (or, more precisely, products of sets together
% with the information of their natural projections to the factors) are final
% objects in
% the category CA,B considered in Example 3.9, for C = Set.
% What is the advantage of viewing products this way? The main advantage
% is that the universal property may be stated in any category, while the
% definition
% of products given in \S{}1.4 only makes sense in Set (and possibly in other
% categories
% where one has a notion of `elements'). We say that a category C has (finite)
% products,
% or is a category `with (finite) products', if for all objects A, B in C the
% category CA,B
% considered in Example 3.9 has final objects. Such a final object consists of the
% data
% of an object of C, usually denoted A \times  B, and of two morphisms A \times  B
% \to  A,
% A \times  B \to  B.
% Note that a `product' from this perspective does not need to `look like a prod-
% uct'. Consider our recurring example of the category obtained from \leq  on Z,
% as
% in Example 3.3. Does this category have products? Objects of this category are
% simply integers a, b \in  Z; call a \times  b for a moment the `categorical'
% product of a and b.
% The universal property written out above becomes, in this case, for all z \in  Z
% such
% that z \leq  a and z \leq  b, we have z \leq  a \times  b.
% This universal problem does have a solution \forall a, b: it is conventionally
% not
% called a \times  b, but rather min(a, b). It is immediate to see that min(a, b)
% satisfies
% the property. Thus this category has products, and in fact we see that the
% product
% in this category amounts to the familiar operation of taking the minimum of two
% integers.
% Thus there is an unexpected connection between `the Cartesian product of two
% sets' and `the minimum of two integers': both are examples of products, taken
% in different categories; they both satisfy `the same' universal property, in
% different
% contexts.

\subsection{Coproducts}
\label{subsec:prelims:coproducts}
% The prefix co- usually indicates that one is `reversing all ar-

% rows'. Just as products are final objects in the categories CA,B obtained by
% consid-
% ering morphisms in C with common source, whose targets are A and B, coproducts
% will be initial objects in the categories23 CA,B of morphisms with common tar-
% get, whose sources are A and B. Dear reader, look away and spell this universal
% property out before we do.
% Here it is. Let A, B be objects of a category C. A coproduct A \sqcup  B of A
% and B will be an object of C, endowed with two morphisms iA : A \to  A \sqcup
% B,
% iB : B \to  A \sqcup  B and satisfying the following universal property: for all
% objects Z
% 23These categories were also considered in Example 3.9; cf. Exercise 3.11.
% and morphisms
% A L L L f
%  L L A
%  L %
% r r r
%  r
%  fB
%  r r 9
%  Z
% B
% there exists a unique morphism \sigma  : A \sqcup  B \to  Z such that the
% diagram
% fA
% A P P P P i
%  P A P
%  '
%  \sigma 
%  $
% n n n
%  n
%  iB
%  n n n
%  A
%  7 \sqcup  B
%  /:
%  Z
% B
% fB
% commutes.
% The symmetry with the universal property of products is hopefully completely
% apparent. We say that a category C has coproducts if this universal problem has
% a
% solution for all pairs of objects A and B.
% Is the reader familiar with any coproduct? Yes!
% Proposition 5.6. The disjoint union is a coproduct in Set.
% Proof. Recall (\S{}1.4) that the disjoint union A \sqcup  B is defined as the
% union of two
% disjoint isomorphic copies A\setminus  , B \setminus  of A, B, respectively; for
% example, we may let
% A\setminus  = {0} \times  A, B \setminus  = {1} \times  B. The functions iA , iB
% are defined by
% iA (a) = (0, a),
%  iB (b) = (1, b),
% where we view these elements as elements of ({0} \times  A) \cup  ({1} \times
% B).
% Now let fA : A \to  Z, fB : B \to  Z be arbitrary morphisms to a common target.
% Define
% \sigma  : A \sqcup  B = ({0} \times  A) \cup  ({1} \times  B) \to  Z
% by
%  ?
% fA (a) if c = (0, a) \in  {0} \times  A,
% \sigma (c) =
% fB (b)
%  if c = (1, b) \in  {1} \times  B.
% This definition makes the relevant diagram commute and is in fact forced upon us
% by this commutativity, proving that \sigma  exists and is unique.
%  ?
% This observation tells us that the category Set has coproducts and further sheds
% considerable light on the mysteries of disjoint unions. For example, there was
% an
% element of arbitrariness in our choice of `a' disjoint union, although different
% choices
% led to isomorphic notions. Now we see why: terminal objects of a category are
% not
% unique in general, although they are unique up to isomorphism (Proposition 5.4);
% there is not a `most beautiful' disjoint union of two sets just as there is not
% a `most
% beautiful' singleton in Set (cf. \S{}5.1).
% Also, an unexpected `symmetry' between products and disjoint unions becomes
% suddenly apparent from the point of view of universal properties.
% The reader is invited to contemplate the notion of coproduct in the other cat-
% egories we have encountered. For example (and probably not surprisingly at this
% point) the category obtained from \leq  on Z does have coproducts: the coproduct
% of
% two objects (i.e., integers) a, b is simply the maximum of a and b.

\subsection*{Exercises}

% 5.1. Prove that a final object in a category C is initial in the opposite
% category Cop
% (cf. Exercise 3.1).
% 5.2. ? Prove that \emptyset  is the unique initial object in Set. [\S{}5.1]
% 5.3. ? Prove that final objects are unique up to isomorphism. [\S{}5.1]
% 5.4. What are initial and final objects in the category of `pointed sets' (Exam-
% ple 3.8)? Are they unique?
% 5.5. ? What are the final objects in the category considered in \S{}5.3?
% [\S{}5.3]
% 5.6. ? Consider the category corresponding to endowing (as in Example 3.3) the
% set Z+ of positive integers with the divisibility relation. Thus there is
% exactly one
% morphism d \to  m in this category if and only if d divides m without remainder;
% there is no morphism between d and m otherwise. Show that this category has
% products and coproducts. What are their `conventional' names? [\S{}VII.5.1]
% 5.7. Redo Exercise 2.9, this time using Proposition 5.4.
% 5.8. Show that in every category C the products A \times  B and B \times  A are
% isomorphic,
% if they exist. (Hint: Observe that they both satisfy the universal property for
% the
% product of A and B; then use Proposition 5.4.)
% 5.9. Let C be a category with products. Find a reasonable candidate for the
% universal property that the product A \times  B \times  C of three objects of C
% ought to
% satisfy, and prove that both (A \times  B) \times  C and A \times  (B \times  C)
% satisfy this universal
% property. Deduce that (A \times  B) \times  C and A \times  (B \times  C) are
% necessarily isomorphic.
% 5.10. Push the envelope a little further still, and define products and
% coproducts
% for families (i.e., indexed sets) of objects of a category.
% Do these exist in Set?
% It is common to denote the product A \times  \dots  \times  A by An .
% ?n times
% 5.11. Let A, resp. B be a set, endowed with an equivalence relation \sim A ,
% resp. \sim B .
% Define a relation \sim  on A \times  B by setting
% (a1 , b1 ) \sim  (a2, b2 ) \iff  a1 \sim A a2 and b1 \sim B b2 .
% (This is immediately seen to be an equivalence relation.)
% \bullet  Use the universal property for quotients (\S{}5.3) to
%          establish that there are func-
% tions (A \times  B)/\sim  \to  A/\sim A, (A \times  B)/\sim  \to  B/\sim B .
% \bullet  Prove that (A\times B)/\sim , with these two functions,
%          satisfies the universal property
% for the product of A/\sim A and B/\sim B .
% \bullet  Conclude (without further work) that (A \times  B)/\sim  \sim 
%  = (A/\sim A ) \times  (B/\sim B ).
% 5.12. \lnot  Define the notions of fibered products and fibered coproducts, as
% terminal ob-
% jects of the categories C\alpha ,\beta  , C\alpha ,\beta  considered in Example
% 3.10 (cf. also Exercise 3.11),
% by stating carefully the corresponding universal properties.
% As it happens, Set has both fibered products and coproducts. Define these
% objects `concretely', in terms of naive set theory. [II.3.9, III.6.10, III.6.11]

%%% Local Variables:
%%% mode: LaTeX
%%% TeX-engine: xetex
%%% TeX-master: "../master"
%%% End:
